\documentclass[11pt,a4paper]{article}
\usepackage[margin=1.8cm]{geometry}
\usepackage{amsmath,amssymb,tabularx,array,needspace,fontspec,enumitem,cancel,booktabs}
\newcommand{\tb}{\cancel{||||}\,}
\setmainfont{FreeSerif}
\newfontfamily\urdufont[Script=Arabic]{FreeSerif}
\TeXXeTstate=1
\newcommand{\urdu}[1]{{\urdufont\beginR #1\endR}}
\setlength{\parindent}{0pt}
\setlength{\parskip}{0pt}
\renewcommand{\arraystretch}{1.15}
\pagestyle{plain}\hyphenpenalty=10000 \exhyphenpenalty=10000

% \lqn{number}{content}
\newcommand{\lqn}[2]{\Needspace{6\baselineskip}\par\medskip\noindent\hangindent=2.0em\hangafter=1\makebox[2.0em][l]{\textbf{#1.}}#2\par}
\begin{document}
\begin{center}
  {\Huge\bfseries The Hope Science Academy}\\[3pt]
  {\small Excellence is our standard not goal}\\[2pt]
  {\small 247-E-1, (Opposite LDA School) Johar Town, Lahore. PH: 0300-8194789, 0322-4157001}
\end{center}

\medskip
\noindent
\begin{tabularx}{\linewidth}{|X|X|X|X|}
\hline
\textbf{Student Name:} & \textbf{Roll Num:} & \textbf{Class:} 11th & \textbf{Paper Code:} \\ \hline
\textbf{Subject Name:} Principles of Accounting & \textbf{Time Allowed:} & \textbf{Total Marks:} 148 & \textbf{Exam Date:} 02-Nov-2026 \\ \hline
\multicolumn{4}{|l|}{\textbf{Exam Syllabus:} FULL BOOK} \\ \hline
\end{tabularx}

\bigskip
\noindent\textbf{Q1. Write detailed answers of the following questions.}\hfill\textbf{1X148=148}

\lqn{1}{Show the effects of following transactions on accounting equation:\\
(a) Mr. Ali started business with cash Rs. 168,000.\\
(b) He purchased goods for cash Rs. 26,000.\\
(c) He purchased furniture costing Rs. 4,000.\\
(d) He sold goods for cash Rs. 6,000 costing Rs. 3,600.\\
(e) He paid salaries Rs. 2,400.\par
}
\lqn{2}{Prepare accounting equation on the basis of the following:\\
(a) Mr. Ferooz started business with Goods Rs. 40,000, Building Rs. 100,000 and Cash Rs. 110,000.\\
(b) Purchased goods on account for Rs. 30,000.\\
(c) Sold goods on account for Rs. 40,000.\\
(d) Cash of Rs. 20,000 deposited into bank.\\
(e) Paid cash to supplier of goods Rs. 20,000.\\
(f) Paid telephone bill Rs. 4,000 by cheque.\par
}
\lqn{3}{Show the effect of following transactions on the accounting equation:\\
(a) Mr. Roshin commenced business with cash Rs. 4200,000.\\
(b) He purchased furniture for Rs. 20,000.\\
(c) He purchased goods worth Rs. 100,000 for cash and goods worth Rs.60,000 on credit.\\
(d) He sold goods to Saleem for Rs. 40,000 costing Rs. 34,000 on credit basis.\\
(e) He withdrew cash Rs. 6,000 for his personal use.\\
(f) He paid rent for building Rs. 4,000.\par
}
\lqn{4}{Show the effect of the following transactions upon the accounting equation:\\
(a) Mr. Babar invested Rs. 2,20,000 in cash.\\
(b) He purchased goods for cash Rs. 44,000 and credit Rs. 36,000.\\
(c) He made additional investment in cash Rs. 40,000.\\
(d) Cash paid to creditors Rs. 20,000.\\
(e) He purchased refrigerator for personal use Rs. 30,000.\\
(f) He paid rent Rs. 4,000 and rent outstanding Rs. 2,000.\par
}
\lqn{5}{Show the effects of the following transactions on the accounting equation of Mr. Qaiser:\\
(a) Commenced business with cash Rs. 2,00,000.\\
(b) Purchased goods for cash Rs. 40,000.\\
(c) Purchased typewriter for cash Rs. 20,000.\\
(d) Rent paid in advance Rs. 17,000.\\
(e) Commission received in advance Rs. 3000.\\
(f) Sold goods for cash Rs. 20,000 costing Rs. 17,000.\par
}
\lqn{6}{Following are the transactions of Wasif \& Co. prepare accounting equation for the month of November 2022:\\
(i) Started business with cash Rs. 200,000.\\
(ii) Bought furniture Rs. 4,000.\\
(iii) Bought machinery Rs. 30,000 paying cash Rs. 10,000 and balance is to be paid later on.\\
(iv) Purchased merchandise Rs. 15,000.\\
(v) Sold merchandise Rs. 8,000 costing Rs. 6,000 to Najam.\\
(vi) Paid electricity bill Rs. 4,000.\\
(vii) Goods Rs. 2,000 used personally by Wasif.\\
(viii) Paid to creditors Rs. 8,000\\
(ix) Received from Najam Rs. 8,000.\\
(x) Received bank interest Rs. 2,000.\par
}
\lqn{7}{Show the effect of the following transaction on accounting equation:\\
(i) Started business with cash Rs. 30,000/.\\
(ii) Purchased goods for cash Rs. 6,000/.\\
(iii) Sold goods to Baqir on account for Rs. 3,000/.\\
(iv) Purchased goods from Imran on account of Rs. 4,000/.\\
(v) Purchased furniture for cash Rs. 5,000/.\\
(vi) Sold goods for Rs. 2,000/. which cost Rs. 1,800/.\\
(vii) Received Rs. 3,000/. from Baqir.\\
(viii) Cash paid to Imran Rs. 4,000/.\\
(ix) Paid rent Rs. 500/.\\
(x) Withdrew cash for personal use Rs. 200/.\par
}
\lqn{8}{Analyze the following transactions:\\
(a) Rizwan commenced business with cash Rs. 1,000,000\\
(b) Purchased goods for cash Rs. 300,000\\
(c) Purchased office furniture for cash Rs. 40,000\\
(d) Paid advertising expenses Rs. 4,000\\
(e) Purchased goods from Adil \& Sons Rs. 28,000\\
(f) Sold goods for cash Rs. 72,000\\
(g) Sold goods to Salman \& Sons Rs. 54,000\\
(h) Received cash from Salman \& Sons Rs. 40,000\\
(i) Paid cash to Adil \& Sons Rs. 20,000\\
(j) Paid salary to Salesman Rs. 3,000\par
}
\lqn{9}{Make the analysis of the following transactions and record them in various accounts by applying the rules of debit and credit.\\
July 01, 2022: Invested Rs. 200,000 cash, building Rs. 150,000 and furniture Rs. 50,000.\\
July 04 20225: Purchased goods for cash Rs. 88,000 subject to trade discount @ 10\%.\\
July 07 2022: Purchased goods from Ali \& Sons for Rs. 10,000.\\
July 12, 2022: Sold goods on account to Zain Rs. 16,000.\\
July 15, 2022: Received cash from Zain Rs. 15,6000, discount allowed Rs. 400.\\
July 16, 2022: Paid Ali \& Sons against account Rs. 9,500, discount received Rs. 500.\\
July 18, 2022: Goods taken away by the owner Rs. 7000.\\
July 21, 2022: Paid salaries Rs. 1,000, Rent Rs. 1,200 and repairing\\
charges for Rs. 800.\\
July 30, 2022: Received commission Rs. 2,400.\\
July 31, 2022: Goods distributed as free samples Rs. 2,000.\par
}
\lqn{10}{Record the following transactions in the journal.\\
August 01, 2022:           Saad started business with cash Rs. 100,000. August 05, 2022:           Purchased goods for cash Rs. 20,000.\\
August 09, 2022:           Sold goods for cash Rs. 24,000.\\
August 15, 2022:           Paid rent Rs. 6,000.\\
August 25, 20225:           Received interest Rs. 21,000.\par
}
\lqn{11}{Record the following transactions in the journal.\\
May 01, 2022:   Commenced business with cash Rs. 500,000.\\
May 03, 2022:   Goods purchased from Tanveer on credit Rs. 20,000. May 08, 2022:   Purchased typewriter from Muzzamal Rs. 20,000. May 15, 2022:   Cash sale to Rahim Rs. 12,000.\\
May 20, 2022:   Paid salary to clerk Rs. 63,000.\par
}
\lqn{12}{Journalise the following transactions in the book of a trader.\\
August 01:        Introduced cash Rs. 296,000 and building Rs. 400,000 as capital.\\
August 04:        Purchased goods on account from Noor Mills limited Rs. 9,000.\\
August 10:        Sold goods to Mobeen \& Co. for Rs. 3,000.\\
August 15:        Goods returned by Mobeen \& Co. for Rs. 250.\\
August 20:        Paid insurance premium Rs. 660.\par
}
\lqn{13}{Journalise the following transactions in the books of Wasim.\\
Jan 01, 2022:    Started business with cash Rs. 800,000.\\
Jan 04, 2022:    Purchased goods for cash Rs. 200,000.\\
Jan 09, 2022:    Purchased motor vehicle from Saleem Auto Rs. 300,000. Jan 14, 2022:    Sold goods to Habib at list price less trade discount 5\% Rs. 100,000.\\
Jan 20, 2022:    Paid travelling expenses Rs. 2,000.\par
}
\lqn{14}{Record the following transactions in journal.\\
Oct 01, 2022:    Umar started business with a capital of Rs. 148,000.\\
Oct 02, 2022:    Purchased goods for Rs. 28,000 on credit from Amjad.\\
Oct 10, 2022:    Sold goods to Aslam for Rs. 18,000.\\
Oct 14, 2022:    Cash given to Amjad Rs. 27,600 and received discount Rs. 400.\\
Oct 19, 2022:    Received cash from Aslam Rs. 17,600 in full settlement of account.\par
}
\lqn{15}{Record the following transactions in journal.\\
Oct 01, 2022:    Started business with a capital Rs. 600,000.\\
Oct 05, 2022:    Purchased goods on credit from Naveed Rs. 60,000. Oct 10, 2022:    Goods returned to Naveed Rs. 2,000.\\
Oct 15, 2022:    Withdrew cash for the payment of tuition fee of owner’s children Rs. 2,000.\\
Oct 20, 2022:    Building purchased for cash Rs. 240,000.\par
}
\lqn{16}{Journalise the following transactions.\\
Mar 01: Commenced business with cash Rs. 500,000.\\
Mar 05: Goods purchased from Moneeb Rs. 80,000.\\
Mar 09: Withdraw goods from business for personal use Rs. 3,000. Mar 14: Purchased machinery Rs. 100,000.\\
Mar 20: Charged depreciation on machinery Rs. 2,000.\par
}
\lqn{17}{Journalise the following transactions in the books of Qaiser.\\
Dec 01, 2022:    Qaiser started business with cash Rs. 400,000. Dec 06, 2022:    Goods purchased from Ansar Rs. 200,000.\\
Dec 10, 2022:    Goods destroyed in accident Rs. 24,000.\\
Dec 14, 2022:    Goods given away as charity Rs. 10,000.\\
Dec 19, 2022:    Cash given away as charity Rs. 2,000.\par
}
\lqn{18}{Record the following transactions in journal.\\
May 01, 2022:   Started business with cash Rs. 12,00,000. May 02, 2022:   Deposited cash into bank Rs. 800,000. May 07, 2022:   Bought goods by cheque Rs. 300,000. May 10, 2022:   Commission paid by cheque Rs. 1,000. May 15, 2022:   Interest received on bank deposit Rs. 500.\par
}
\lqn{19}{Record the following transactions in the journal.\\
August 01, 2022:           Saad started business with cash Rs. 100,000. August 05, 2022:           Purchase goods for cash Rs. 20,000 and credit for Rs. 10,000 from Ali.\\
August 07, 2022:           Purchased furniture for cash Rs. 8,000.\\
August 09, 2022:           Sold goods for cash Rs. 24,000 and on account Rs. 12,000 to Aslam.\\
August 15, 2022:           Ali account settled receiving Rs. 400/- as discount. August 25, 2022:           Aslam account settled allowing Rs. 400/- as discount.\par
}
\lqn{20}{Journalize the following business transactions of Faisal Enterprises:\\
(i) Faisal started business with cash Rs, 50,000/-\\
(ii) Purchased building for cash Rs. 10,000/.\\
(iii) Bought merchandize on amount Rs. 15,000/.\\
(iv) Purchased furniture for cash Rs. 5,000/.\\
(v) Sold goods for cash Rs. 10.000/.\\
(vi) Paid cash to creditors Rs. 7,000/. for goods purchased on account.\\
(vii) Sold goods for customers on account Rs. 6,000/.\\
(viii) Withdraw cash Rs. 2,000/. for personal use.\\
(ix) received cash Rs. 3,000/. from customers for goods sold to them on account.\\
(x) Paid salary to employee Rs. 3,500/.\par
}
\lqn{21}{Journalize the following transactions:\\
(i) Yasir started business with cash Rs. 50,000/.\\
(ii) Purchased furniture for cash Rs. 10,000/.\\
(iii) Purchased goods for cash Rs. 3,000/.\\
(iv) Purchased goods from Ali on account Rs. 5,000/.\\
(v) Sold goods for cash Rs. 8,000/.\\
(vi) Sold goods to Abid on account Rs. 4,000/.\\
(vii) Paid cash to Ali Rs. 5,000/.\\
(viii) Received cash from Abid Rs. 4,000/.\\
(ix) paid rent of shop Rs. 3,000/.\\
(x) Yasir withdrew cash Rs. 2,000 for his personal use.\par
}
\lqn{22}{Record the following transactions in journal.\\
July 01, 2022:    Umar started business with a capital of Rs. 148,000 of which Rs. 68,000 was paid into bank.\\
July 02, 2022:    Purchased goods for cash Rs. 28,000 and Rs. 28,000 on credit from Amjad.\\
July 04, 2022:    Purchased office typewriter for Rs. 22,000 from Muzammal.\\
July 06, 2022:    Sold goods for cash Rs. 8,000 of which Rs. 2,000 was paid into bank.\\
July 10, 2022:    Sold goods for Aslam for cash Rs. 8,000 and to Alam for Rs. 18,000.\\
July 14, 2022:    Cheque given to Amjad Rs. 27,600 and received discount Rs. 400.\\
July 17, 2022:    Amount withdrawn from bank for personal use Rs. 3,000 and for business use Rs. 6,000.\\
July 19, 2022:    Received a cheque from Alam Rs. 17,600 in full settlement of account.\\
July 30, 2022:    Purchased stationary for office Rs. 3,200.\\
July 31, 2022:    Paid rent Rs. 2,400, salaries Rs. 6,800 by cheque.\par
}
\lqn{23}{Journalise the following transactions in the books of a trader.\\
August 01, 2022:           Introduced cash as capital Rs. 2,96,000.\\
August 01, 2022:           Purchased furniture for cash Rs. 60,000.\\
August 03, 2022:           Cash purchases for Rs. 32,000.\\
August 04, 2022:           Purchased goods on account from Noor Mills limited Rs. 9,000.\\
August 09, 2022:           Goods returned to Noor Mills limited for Rs. 1000. August 12, 2022:           Cash sales for Rs. 19,800.\\
August 18, 2022:           Paid to Noor Mills Ltd. Rs. 7,600 in full settlement of account.\\
August 21, 2022:           Sold goods to Mobeen \& Co. for Rs. 3,000.\\
August 22, 2022:           Loss by accident while goods were in transit Rs. 150.\\
August 24, 2022:           Goods returned by Mobeen \& Co. for Rs. 250. August 25, 2022:           Paid tuition fee of owner children Rs. 600.\\
August 26, 2022:           Received discount for Rs. 100 from Mobeen \& Co. and settled his account.\\
August 30, 2022:           Paid insurance premium Rs. 660.\par
}
\lqn{24}{Journalise the above transactions and post them into the ledger.\\
Mr. Saeed started his business with cash Rs. 130,000 and Building Rs. 240,000 on 1\textsuperscript{st}\\
January, 2022. His transactions during the month were as follows: Jan 2, 2022:      Purchased goods for cash Rs. 60,000.\\
Jan 5, 2022:      Sold goods for cash Rs. 44,000.\\
Jan 10, 2022:    Bought an iron safe for Rs. 8,000.\par
}
\lqn{25}{\\
Record the following transactions in Journal and post them in ledger and prepare a trial balance:\\
May 1, 2022:     Business started with cash Rs. 200,000.\\
May 4, 2022:     Goods purchased for cash from Anwar Rs. 60,000. May 7, 2022:     Sold goods to Rahim Bros Rs. 20,000.\\
May 14, 2022:   Paid salaries Rs. 10,000.\\
May 20, 2022:   Commission received Rs. 900.\\
May 25, 2022:   Cash received from Rahim Bros in full settlement Rs. 19,800.\par
}
\lqn{26}{\\
Record the following transactions in the journal, post them in the ledger and prepare a trial balance.\\
March 01, 2022: Mr. Ajmal started business with Rs. 160,000 cash, Machinery worth Rs. 20,000 and Stock worth Rs. 80,000.\\
March 04, 2022: Deposited into bank Rs. 90,000.\\
March 09, 2022: Bought goods on Credit from Noor Sons Rs. 40,000.\\
March 14, 2022: Sold goods to Imran for Rs. 24,000.\\
March 25, 2022: Issued a cheque to Noor Sons of Rs. 39,600 in full settlement of his account.\\
March 30, 2022: Received cash from Imran Rs. 23,500 in full settlement of his account.\par
}
\lqn{27}{\\
Record the following transactions in the journal, post in the ledger and prepare a trail balance.\\
March 01, 2022: Mr. Ajmal started business with Rs. 160,000 cash, Machinery worth Rs. 20,000 and Stock worth Rs. 80,000.\\
March 02, 2022: Deposited into bank Rs. 905,000.\\
March 03, 2022: Bought goods on Credit from Noor Sons Rs. 40,000.\\
March 06, 2022: Sold goods to Imran for Rs. 24,000.\\
March 08, 2022: Paid Travelling expenses Rs. 1,200.\\
March 13, 2022: Issued a cheque to Noor Sons of Rs. 99,600 in full settlement of his account.\\
March 16, 2022: Received cash from Imran Rs. 23,500 in full settlement of his account.\\
March 18, 2022: Bank charges made by cheque Rs. 150.\\
March 22, 2022: Goods sold for cash Rs. 12,000.\\
March 24, 2022: Paid rent of building by cheque Rs.600.\\
March 30, 2022: Interest received Rs. 800.\\
March 31, 2022: Paid carriage Rs. 1,600.\par
}
\lqn{28}{Give journal entries in the books of Afan and Azan\\
Afan sells goods to Azan for Rs. 20,000 on 1st April 2022, he draws on him a bill for three months which Azan accepts, at maturity the bill is honored.\par
}
\lqn{29}{Show the entries in the books of A, B \& Bank.\\
A Draws a bill for Rs. 4,000 on B, who accepts and returns it to A on the same date. A discounted the bill with his banker for Rs. 3,900, the bill is finally honored by B.\par
}
\lqn{30}{Give journal entries in the books of Arif, Hanif and Haseeb.\\
Arif sold goods to Hanif for Rs. 40,000 on credit on 1\textsuperscript{st}\\
March 2022. Arif drew and Hanif accepted a bill for Rs. 40,000 at three months on the same date. Haseeb was creditor of Arif for Rs. 40,000. Arif settled his account by endorsing the bill to Haseeb on March 3, 2022. On maturity, Hanif honored his acceptance.\par
}
\lqn{31}{Give journal entries in the books of A and Bank.\\
A sold goods to B for Rs. 9,000 on credit on 1\textsuperscript{st}\\
Jan, 2022. A drew a bill on B for Rs. 9,000 on the same date at ‘3’ months. On 2\textsuperscript{nd}\\
Jan, 2022, A sent the bill to the bank for collection. On the due date B met his acceptance and bank received Rs. 40 as collection charges.\par
}
\lqn{32}{Give entries in the books of A and B.\\
On 1\textsuperscript{st}\\
September, 2022 A owes B a sum of Rs. 6,000. He gives him two acceptances one for Rs. 4,400 at 2 months and 2\textsuperscript{nd}\\
for Rs. 1,600 at 3 months. B endorses the 1\textsuperscript{st}\\
bill in the favor of his creditor C and sent the 2\textsuperscript{nd}\\
bill to his bank for collection. Both the bills are honored by ‘A’ on the due date. The bank charges Rs. 40 as collection charges.\par
}
\lqn{33}{B \& Co. sells goods to P for Rs. 12,000 and draws a bill for four months, P accepts the bill and returns it to B \& Co. Show what entries should be passed in the books of B \& Co.\\
(a) If they retained the bill till the due date and then realized on maturity.\\
(b) If they discounted the bill with their banker for Rs. 11,760.\\
(c) If they endorsed the bill to their creditor M \& Co. in settlement of debt.\\
(d) If they sent the bill to the Bank for collection.\par
}
\lqn{34}{Rehman sells goods to Imran for Rs. 60,000 and draws a bill for two months. Imran accepts the bill and returns it to Rehman. Pass the journal entries in the books of Rehman and Imran. If it is dishonored at maturity.\par
}
\lqn{35}{Give journal entries in the books of X, Y and Bank.\\
X draws a bill for Rs. 20,000 on Y, who accepts and returns it to X on the same date for ‘3’ months. X discounted the bill with his banker at 9\% p.a. On the due date the bill was dishonored by Y.\par
}
\lqn{36}{Give entries in the books of A, B and C.\\
On 1\textsuperscript{st}\\
September, 2022 A owes B a sum of Rs. 6,000. He gives him acceptance for Rs. 6,000 at 2 months. B endorses the bill in the favor of his creditor C. Bill is dishonored by ‘A’ on the due date.\par
}
\lqn{37}{Give journal entries in the books of B, A and Bank.\\
On 1\textsuperscript{st}\\
January, 2022 A owes B a sum of Rs. 4,000. He gives him acceptance for Rs. 4,000 at ‘3’ months. B sent the bill to his bank for collection. Bill dishonored on due date by A and bank paid Rs. 40 as noting charges.\par
}
\lqn{38}{B \& Co. sells goods to P for Rs. 12,000 and draws a bill for four months, P accepts the bill and returns it to B \& Co. Show what entries should be passed in the books of B \& Co. If the bill is dishonored in the following cases:\\
(a) If they retained the bill till the due date.\\
(b) If they dishonored the bill with their banker for Rs. 11,760.\\
(c) If they endorsed the bill to their creditor M \& Co. in settlement of debt.\\
(d) If they sent the bill to the Bank for collection.\par
}
\lqn{39}{Make the necessary journal entries to record these transactions in A’s journal.\\
A receives a bill of exchange from B for Rs. 200 and discounts with his bankers for Rs. 192. On the bill becoming due; B requests A to renew it which A does, adding Rs. 4 to the new bill for interest. On the due date of the second bill ‘B’ met his acceptance.\par
}
\lqn{40}{On 1\textsuperscript{st}  January, 2022 Sohail sold goods to Zaheer for Rs. 10,000 and drew a bill for 4 months on him. Zaheer accepted the bill and returned it to Sohail. Sohail got it discounted by his bank at 15\% p.a. On Maturity Zaheer was not in a position to meet his acceptance. He approached Sohail and requested him to cancel the old bill and draw on him a new bill for the amount of old bill plus interest at the rate of 12\% p.a. at 3 months. Sohail agreed to the proposal. On the due date of the new bill, Zaheer honored his acceptance. Give journal entries in books of Sohail and Zaheer.\par
}
\lqn{41}{Pass the journal entries in the books of B and A.\\
A owes B a sum of Rs. 3000, on 1\textsuperscript{st}\\
January, 2022, he accepts a bill for the amount for 3 months drawn on him by B. B gets it discounted with his banker for Rs. 2,900. On the due date the bill is dishonored and the bank, pays Rs. 40 as noting charges. A pays Rs. 800 in cash and accepts another bill for the balance together with interest of Rs. 60, for two months. On the due date the bill is honored by the acceptor A.\par
}
\lqn{42}{Pass the journal entries in the books of A, B and Bank.\\
A drew on B a bill for Rs. 30,000 at 3 months. A got the bill discounted by his bank at 12\% p.a. On maturity the bill is dishonored and bank pays Rs. 60 for noting charges.\par
}
\lqn{43}{Show the entries in the books of Hassan.\\
On 1\textsuperscript{st}\\
January 2022, Hassan sells goods to Taimoor for Rs. 10,000. Hassan\\
draws a bill on Taimoor for Rs. 10,000 for two months. The bill is discounted at bank on 4\textsuperscript{th}\\
January for Rs. 9,700. The bill is dishonored and the bank pays Rs. 40 for noting charges. On Taimoor’s request, Hassan draws a second bill on Taimoor for Rs. 10,100 including Rs. 60 as interest at two months. On maturity, the second bill is honored.\par
}
\lqn{44}{Journalise the following transactions in the books of Asif:\\
(a) \\
Saleem’s acceptance to Asif for Rs. 4,000 which Asif endorsed in favor of Rauf was dishonored. Rauf paid Rs. 40 as noting charges. Asif paid Rs. 40 to Rauf by a cheque and received a new acceptance from Saleem for the amount due plus interest Rs. 80.\\
(b) Asif renewed his acceptance to Amir for Rs. 2,000 by a cheque of Rs. 800 and a new bill at 3 months at 10\% p.a. interest.\par
}
\lqn{45}{Journalise the following transactions in the books of Asif:\\
(a) Farid’s acceptance to Asif for Rs. 6,000 renewed for two months at 12\% p.a. interest.\\
(b) Asif’s acceptance to Haseeb for Rs. 10,000 was got retired one month before its maturity under a rebate of 12\% p.a.\\
(c) Asif’s acceptance to Ansar for Rs. 2,000 was discharged by Javed’s acceptance to Asif for a similar amount.\\
(d) Asif’s acceptance to Farooq for Rs. 10,000 was discharged by paying Rs. 4,000 in cash and accepting a new bill for the balance for two months, interest at 12\% p.a. being paid in cash.\par
}
\lqn{46}{Pass the entries in the books of Salman and Rizwan and also give Rizwan’s ledger account.\\
Salman sold goods to Rizwan for Rs. 100,000 on 1\textsuperscript{st}\\
July 2022 and drew on him a bill for the amount payable after four months. Rizwan accepted the bill. Salman got it discounted by his bank at 12\% p.a. On maturity Rizwan was in a position to pay only Rs. 40,000 in cash, therefore, he approached Salman for the renewal of the bill, Salman agreed and drew a new bill for Rs. 60,000 along with an interest at 15\% p.a. for 3 months. Before the due date of the new bill Rizwan became insolvent and a first and final divided of 50\% was received from his estate.\par
}
\lqn{47}{Journalise the transactions in the books of both parties.\\
Imran sells goods to Farooq for Rs. 16,000 and draws a bill for 4 months on him. Farooq accepts the bill and returns it to Imran. After three months Farooq gets the bill retired under a rebate of 12\% p.a\par
}
\lqn{48}{Give journal entries in books of A and B.\\
For mutual accommodation, A draws a bill on B for Rs. 10,000 at two months which is accepted by B. On the same date B draws a bill on A for a similar amount at two months. Both the parties get their bills discounted by their respective banks at 12\% p.a. On the due date both the bills are honoured.\par
}
\lqn{49}{Pass journal entries in the books of Waqas and Kashif.\\
Waqas and Kashif are in need of funds. For their mutual accommodation Waqas draws a bill on Kashif for Rs. 12,000 at two months, which Kashif accepts, Waqas gets the bill discounted by his bank at 12\% p.a. and sends half of the proceeds to Kashif. Before he due date Waqas sends the remaining amount due to Kashif and Kashif pays the bill on the due date.\par
}
\lqn{50}{Pass the journal entries in the books of Amin.\\
On 1\textsuperscript{st}\\
October 2014, Amin drew a bill for Rs. 36,000 on Bilal payable after two months for their mutual accommodation to the extent of 3/4 and 1/4 respectively. The bill was duly accepted and returned by Bilal. Amin got it discounted by his bank at 10\% p.a. and remitted 1/4 of the proceeds to Bilal, before the due date Amin sent him the balance but Bilal dishonored the bill on the due date and Amin had to pay his banker the amount of the bill along with noting charges Rs. 60.\par
}
\lqn{51}{Give journal entries and Aqeel’s account in the books of Fida.\\
For their mutual accommodation, Fida accepted a bill on 1\textsuperscript{st}\\
March 2022, drawn on him by Aqeel for Rs. 12,000 at three months. The bill was got discounted at 10\% p.a. and the proceeds were shared equally.\\
On 1\textsuperscript{st}\\
April 2022, Fida drew a bill for Rs. 18,000 on Aqeel at three months for the same purpose, Aqeel accepted the bill. Fida got the bill discounted at 10\% p.a. and the proceeds were shared as 2/3 to Fida and 1/3 to Aqeel. Before the due date of the first bill. Fida sent a cheque to Aqeel in full settlement of his account. On 15\textsuperscript{th}\\
June 2022, Aqeel became insolvent and his estate paid a divided of 50 paisa in a rupee.\par
}
\lqn{52}{Pass journal entries and show account of Ferooz in the books of Roshin. Ferooz for mutual accommodation draws a bill for Rs. 15,000 on Roshin at three months. Ferooz gets the bill discounted by his banker for Rs. 14,625 and remits Rs. 4,875 to Roshin. On maturity Ferooz is not able to send the amount due to Roshin, to enable him to meet the bill. He, however, accepts a bill for Rs. 18,750 which is discounted by Roshin for Rs. 17,625. Roshin meets his acceptance and remits Rs. 900 to Ferooz. Before the due date of the bill Ferooz becomes insolvent and a dividend of 60 paisa in the rupee is received from his estate.\par
}
\lqn{53}{Write up the purchases Journal (Book) of Mr. Asif from the following particulars and show their posting into the ledger.\\
Jan 01, 2022:    Purchased from Bashir Plastic Co. Ltd. 12 dozens of plastic tubes @ Rs. 140 per dozen.\\
Jan 02, 2022:    Bought of Imran \& Sons 20 table glasses @ Rs. 30 each. Jan 10, 2022:    Bought of Rafiq Ltd. 30 tables @ Rs. 120 each and 5 chairs @ Rs. 200 each.\\
Jan 18, 2022:    Bought of Riaz \& Sons 50 shirts @ Rs. 100 each less 10\% trade discount.\par
}
\lqn{54}{From the following particulars write up the sales Journal of Raza computers and show their posting into the ledger.\\
April 02, 2022:   Sold to Ali \& Bros on credit 14 hard disks @ Rs. 2000 each.\\
April 08, 2022:   Sold to Asim \& Sons 20 modems @ Rs. 1100 each.\\
April 15, 2022:   Sold to Saad computers 14 processors @ Rs. 2,500 each. April 22, 2022:   Sold to Tariq \& Sons on credit 12 monitors @ Rs. 4,000 each, less trade discount 10\%.\par
}
\lqn{55}{Enter the following transactions in the Sales Returns Journal (Returns Inwards Book) and post them to the ledger accounts.\\
April 03, 2022:   Received back from Sohail \& Co. 10 tables @ 100 per table not according to specification.\\
April 12, 2022:   Returned by K. Khalid an almirah valued at Rs. 600 for damaged in transit.\\
April 18, 2022:   Received back from M. Ashraf 12 chairs @ Rs. 50 each not according to order.\\
April 20, 2022:   Returned by Waseem \& Bros a show case valued at Rs. 1,000 for inferior quality.\par
}
\lqn{56}{From the following particulars write the purchases Returns Book and show the posting to the ledger.\\
Jan 10, 2022:    Returned goods to Sadiq \& Sons. Being of inferior quality valued at Rs. 500.\\
Jan 15, 2022:    Returned goods to Lahore stores being not according to sample valued at Rs. 1,000.\\
Jan 18, 2022:    Returned goods to Danish \& Co. being damaged in transit valued at Rs. 800.\\
Jan 25, 2022:    Returned goods to Usman stores. Being not according to specification, valued at Rs. 400.\par
}
\lqn{57}{Enter the following transactions in the relevant subsidiary books (special journals) of a trader and post them in the ledger accounts.\\
(a) Purchased goods from Farids Bros. Rs. 5000.\\
(b) Sold goods to X on credit Rs. 3000.\\
(c) Goods purchased from Farid \& Bros. valued at Rs. 500 now returned to them being not according to sample.\\
(d) Goods worth Rs. 500 sold to X now received back from him as these were found, unsuitable.\\
(e) Purchased, from Noor \& Co. 100kg tea @ 30 per kg.\\
(f) Sold to Naeem \& Sons 30kg. Tea @ Rs. 40 per kg.\\
(g) Returned by Naeem \& Sons 10kg, tea being not according to sample.\\
(h) Returned 10kg tea to Noor \& Co. being damaged.\par
}
\lqn{58}{Enter the following transactions in a Single Column Cash book of Irfan.\\
(a) Cash in hand Rs. 50,000.\\
(b) Purchased goods for cash Rs. 4,000.\\
(c) Purchased a typewriter for business Rs. 6,000.\\
(d) Sold goods to Waseem for Rs. 2,400 @ 10\% trade discount and cash received.\\
(e) Paid to Mr. Rehan Rs. 1,500.\\
(f) Cash Sales Rs. 1000.\\
(g) Received from Ahmed \& Co. Rs. 2,700.\\
(h) Purchased stationary for Cash Rs. 1,600.\\
(i) Received from Zaheer \& Co. Rs. 4,800.\\
(j) Paid office rent Rs. 1000.\\
(k) Paid salaries Rs. 1,500.\par
}
\lqn{59}{Enter the following transactions in a Double Column Cash Book (with cash and bank column) of Khurram \& Co. and find out the balance on 07-01-2022.\\
Jan 01, 2022:    Cash in hand Rs. 100,000\\
Jan 01, 2022:    Cash at Bank Rs. 14,00,000.\\
Jan 02, 2022:    Paid salaries Rs. 80,000.\\
Jan 02, 2022:    Paid rent by cheque Rs. 60,000.\\
Jan 04, 2022:    Bought goods by cheque Rs. 860,000.\\
Jan 05, 2022:    Sold goods for cash Rs. 150,000.\\
Jan 06, 2022:    Bought machinery by cheque Rs. 250,000.\\
Jan 06, 2022:    Deposited into bank Rs. 14,000.\\
Jan 07, 2022:    Withdrew from bank for business use Rs. 90,000.\\
Jan 07, 2022:    Cash sales Rs. 180,000.\\
Jan 07, 2022:    Withdrew cash from business for personal use Rs. 80,000. Jan 07, 2022:    Received bank interest Rs. 6,000.\par
}
\lqn{60}{Record the following transactions into the Cash Book (Double Column) of Fazal Corp and find out the balances at the end of the month.\\
Sep 1, 2022:      Cash balance Rs.135,400.\\
Sep 2, 2022:      Opened a current A/c with HBL and deposited cash therein Rs. 100,000.\\
Sep 5, 2022:      Purchased furniture and paid by cheque Rs. 9,000.\\
Sep 6, 2022:      Cash sales Rs. 2,000.\\
Sep 10, 2022:    Received from Shoiab cash Rs. 1000 and a cheque Rs. 6,000 and deposited the cheque into the Bank.\\
Sep 15, 2022:    Purchased goods and paid by cheque Rs. 11,600.\\
Sep 20, 2022:    Received cash as commission Rs. 1,200.\\
Sep 23, 2022:    Received cash in full settlement from Khalid and sent the cash into bank Rs. 1000.\\
Sep 29, 2022:    Withdrawn from bank for office use Rs. 10,000.\\
Sep 30, 2022:    Paid salary Rs. 15,000.\\
Sep 30, 2022:    Paid rent Rs. 800.\par
}
\lqn{61}{From the under mentioned information write up a Double Column Cash Book and strike out the balance at the end of the month.\\
March 01, 2022: Cash in hand Rs. 60,000.\\
March 01, 2022: Cash at bank Rs. 80,000.\\
March 05, 2022: Received two cheques: (a) from A. Rs. 20,000 and (b) from B. Rs. 10,000.\\
March 12, 2022: Cash purchases Rs. 16,000 less 5\% Trade discount. March 20, 2022: Cash sales Rs. 4,000 and banked the same.\\
March 22, 2022: Withdrawn from the bank Rs. 16,000 and paid wages Rs. 6,000 and rent Rs. 10,000.\\
March 23, 2022: Received from X cash Rs. 6,000 and cheque Rs. 10,000 both deposited into bank.\\
March 25, 2022: Received cheque for sale of old machinery and banked Rs. 16,000\\
March 28, 2022: Paid to Y cash Rs. 4,000 and cheque Rs. 6,000.\\
March 30, 2022: Bank charged for commission Rs. 1,000\par
}
\lqn{62}{From the following particulars prepare a Treble Column Cash Book.\\
Dec 01. 2022:    Cash in hand Rs. 3,000 and at bank Rs. 140,000.\\
Dec 02. 2022:    Received from Rameez Rs.19,8000. Discount allowed Rs. 200.\\
Dec 03. 2022:    Deposited into the bank Rs. 10,000.\\
Dec 04. 2022:    Purchased goods for Rs. 50,000 and paid by cheque.\\
Dec 05. 2022:    Sold goods to Amir on credit Rs. 14,000.\\
Dec 07. 2022:    Received from Amir a cheque for Rs. 14,000.\\
Dec 08. 2022:    Sold goods for cash Rs. 14,000.\\
Dec 09. 2022:    Paid Saleem Rs. 19.600 by cheque, discount received Rs. 400.\\
Dec 11. 2022:    Paid telephone charges Rs. 1,200.\\
Dec 13. 2022:    Purchased goods for cash Rs. 8,000.\\
Dec 19. 2022:    Withdrew from the bank for business use Rs. 20,000.\\
Dec 25. 2022:    Paid rent by cheque Rs. 1,400.\\
Dec 31. 2022:    Stationary purchased for cash Rs. 300.\par
}
\lqn{63}{Prepare three column cash book.\\
January 01, 2022:          Javaid started business with cash Rs. 192,000. January 02, 2022:          Deposited with bank Rs. 52,000.\\
January 04, 2022:          Purchased goods from Hamid for Rs. 3,200 issued cheque for payment.\\
January 05, 2022:          Purchased Iron safe for Rs. 2,400, paying cash Rs. 1,200, issued cheque for the balance.\\
January 06, 2022:          Purchased goods on credit from Ali for Rs. 4,400. January 07, 2022:          Sold goods on account to Salim for Rs. 4,400. January 10, 2022:          Defective goods returned by Salim Rs. 200. January 11, 2022:          Received a cheque from Salim for Rs. 4,200. January 12, 2022:          Paid Salim’s cheque into bank.\\
January 28, 2022:          Goods sold to Kashif for cash Rs. 6,600.\\
January 29, 2022:          Cheque issued to Ali for Rs. 4,300 and discount received Rs. 100.\\
January 30, 2022:          Salaries paid Rs. 6,600 by cheque.\par
}
\lqn{64}{Prepare a Petty Cash Book on the Imprest system from the following transactions.\\
Jan 01, 2022:    Received Rs. 1,000 from the Chief Chaier for Petty Cash. Jan 02, 2022:    Paid bus fare Rs. 20.\\
Jan 02, 2022:    Paid cartage Rs. 40.\\
Jan 03, 2022:    Paid for postage and telegrams Rs. 90.\\
Jan 03, 2022:    Paid Wages to casual laborers Rs. 100.\\
Jan 04, 2022:    Paid for stationary Rs. 130.\\
Jan 04, 2022:    Paid bus fare Rs. 70.\\
Jan 05, 2022:    Paid for repair of chairs Rs. 200.\\
Jan 07, 2022:    Paid cartage Rs. 170.\\
Jan 07, 2022:    Paid for postage and telegrams Rs. 120.\par
}
\lqn{65}{From the following particulars prepare a Bank reconciliation statement of Sarmad \& Company as on 31\textsuperscript{st}  December 2022.\\
i) Balance as per cash book Rs. 25,000.\\
ii) Cheques issued amounting to Rs. 9,000 not presented, for payment. iii) Cheques deposited but not cleared Rs. 6,220.\\
iv) Bankers had wrongly debited the firm’s account with Rs. 1,424.\\
v) Bank charges debited in pass book Rs. 70.\par
}
\lqn{66}{Prepare a Bank Reconciliation statement with the help of following information as on 31\textsuperscript{st}  December 2022.\\
(a) Balance as per cash book Rs. 60,000.\\
(b) Profit on account credited by the bank Rs. 3,000.\\
(c) Bank paid Rs. 10,000 as rent on the behalf of account-holder.\\
(d) Bank collected dividend on shares Rs. 5,000.\\
(e) Cheques paid into bank but not collected by bank Rs. 8,000.\\
(f) Cheques issued but not presented to the bank for payment Rs. 6,000.\par
}
\lqn{67}{Prepare a bank Reconciliation statement as on 31\textsuperscript{st} March 2022. According to the cash book of Sarfraz Ltd. Company has a balance at the bank Rs. 7,600 on 31\textsuperscript{st}\\
March, 2022 but this did not agree with the bank statement on the same date. An investigation into the difference gives the following information:\\
(i) Trade subscription Rs. 1000 paid by the bank on 31\textsuperscript{st}\\
March 2022 has not been entered in cash book.\\
(ii) A cheque paid for advertising on 10\textsuperscript{th}\\
March for Rs. 394 has not been entered in cash book.\\
(iii) Cheques for Rs. 2,200 sent to creditors on 30\textsuperscript{th}\\
March 2022 were not paid by the book until 5\textsuperscript{th}\\
April 2022.\\
(iv) Cheques received from customers amounting to Rs. 3,360 were paid into bank on 28\textsuperscript{th}\\
March, but were not credited by the bank.\\
(v) Interest on investment collected by bank Rs. 600, not yet recorded in cash book.\par
}
\lqn{68}{Prepare a Bank Reconciliation Statement as on 31\textsuperscript{st} December 2022 from the following particulars:\\
(i) Bank balance as per the Cash Book on 31\textsuperscript{st}\\
Dec. 2022. Rs. 47,020\\
(ii) Cheques for Rs. 6,000 deposited for collection but not credited by the bank prior to 31\textsuperscript{st}\\
December 2022.\\
(iii) Cheques amounting to Rs. 8,000 were issued on 28\textsuperscript{th}\\
December, were not cashed up to 31\textsuperscript{st}\\
December 2022.\\
(iv) A wrong debit of Rs. 800 appeared in Pass Book.\\
(v) Bank charges Rs. 20 appeared in the Pass Book but not in the Cash Book.\par
}
\lqn{69}{Prepare a bank Reconciliation Statement.\\
The balance as per cash book for Aslam Traders on 31\textsuperscript{st}\\
January 2022 was Rs. 38,740:\\
(i) Cheques amounting to Rs. 5,640 issued to a creditor and entered in cash book were not presented for payment.\\
(ii) A cheque for Rs. 3,500 received from Universal Ltd. and paid into bank was dishonoured.\\
(iii) An out-station cheque amounting to Rs. 1,150 entered in cash book as paid into bank but has not been collected, by bank.\\
(iv) Bank charges Rs. 50 have not been entered in cash book.\\
(v) Interest on Govt. Securities Rs. 600 collected and credited in pass book.\par
}
\lqn{70}{From the following particulars prepare a Bank reconciliation statement as at 31\textsuperscript{st}  December 2022 of Salman:\\
(i) Overdrawn as per cash book Rs. 10.200\\
(ii) Interest on Bank overdraft debited in pass book only Rs. 230\\
(iii) Cheques issued but not cashed from bank Rs. 1,800\\
(iv) Cheques paid in but not cleared by bank Rs. 2,750\\
(v) B/R discounted with the bank in November but dishonoured 31\textsuperscript{st} December Rs. 630\par
}
\lqn{71}{From the following particulars, prepare a Bank reconciliation statement of Mr. Khubaib as on 31\textsuperscript{st}  December, 2022:\\
(a) Balance as per pass book Rs. 10,868\\
(b) Cheque issued but not presented Rs. 2,120\\
(c) Cheque deposited but not credited by bank Rs. 1,380\\
(d) Interest charged by bank Rs. 72\\
(e) Interest on Govt. securities Rs. 140\par
}
\lqn{72}{From the following particulars ascertain the bank balance as per the cash book of Naveed \& Co. as on 31\textsuperscript{st}  March 2022:\\
(a) Bank balance as per the pass book on 31\textsuperscript{st} March 2022 was Rs. 160,000\\
(b) Cheques paid into Bank on 26\textsuperscript{th} March 2022 but not collected prior to 31\textsuperscript{st} March 2022 amount to Rs.40,000\\
(c) Interest on deposited of Rs. 4,000 was credited in the pass book but not recorded in the Cash Book till 31\textsuperscript{st}\\
 March 2022\\
(d) A customer paid into the bank Rs. 30,000 directly as appeared in the pass book but not in the Cash Book\\
(e) Dividend money of Rs. 16,000 on the shared held by Naveed \& Co. was received directly by the bank on 29\textsuperscript{th}\\
 March 2022 but the intimation was sent to them no 2\textsuperscript{nd}  April 2022\par
}
\lqn{73}{Prepare a Bank Reconciliation Statement from the following particulars:\\
(a) Overdraft balance as per the Pass Book (Dr.) on 31-12-2022 Rs.10,000\\
(b) Interest on overdraft charged by the bank for half year Rs. 550 debited in the Pass Bank only\\
(c) Two cheques for Rs. 15,000 and Rs. 10,000 paid in on 30\textsuperscript{th}\\
 December 2022 were not credited by the bank\\
(d) Two cheques for Rs. 2,100 and Rs. 8,000 issued in December 2022 had not been presented at the bank during the month\\
(e) A cheque for Rs. 4,000 entered in the cash book (debit side) was not sent to the bank for collection.\par
}
\lqn{74}{A trader has two Bank Account -{}- A/c No. 1 and A/c No. 2. The following particulars relating to A/c No. 1 are available on 31\textsuperscript{st} March 2022. Prepare a bank reconciliation statement as on 31\textsuperscript{st} March 2022:\\
(1) Pass Book balance (overdrawn) Rs. 60,000\\
(2) Cheques drawn prior to 31\textsuperscript{st}\\
 March 2022 but not presented as yet Rs. 16,000\\
(3) Interest debited by the bank but not entered in the Cash Book as yet Rs. 1,600\\
(4) Transfer from A/c No. 2 to A/c No. 1 recorded by the Bank on 31\textsuperscript{st}\\
 March but not entered in the Cash Book Rs. 12,000\\
(5) Bank Charges debited by the bank but not recorded in the Cash Book as yet Rs. 40\par
}
\lqn{75}{On 31\textsuperscript{st} January 2022, a trader's cash book showed an overdraft of Rs. 4,000 at the bank. On reference to the book the following items were noted, prepare bank reconciliation statement:\\
(a) Interest amounting to Rs. 300 had been credited by the Bank in the\\
pass book only\\
(b) Bank charges amounting to Rs. 50 had been debited by the bank in the pass book only.\\
(c) An outstation cheque for Rs. 310 banked on 29\textsuperscript{th}\\
 January but not yet been collected.\\
(d) Cheques drawn totaling Rs. 190 had not been presented at the Bank for payment\\
(e) The credit side of the cash book was undercast by Rs. 200\par
}
\lqn{76}{On 30\textsuperscript{th} June 2022 the Cash Book of Mr. Zubair showed an overdraft balance of Rs. 32,000. This balance did not agree with the balance as per the pass book. On comparing the two books, the following discrepancies were noted. Prepare a bank reconciliation statement with the help of revised cash book:\\
1) Cheques for Rs. 8,000 paid into the bank for collection on 28\textsuperscript{th}\\
 June 2022 have not yet been collected\\
(2) Cheques for Rs.6,000 issued on 26\textsuperscript{th}\\
 June 2022 have not yet been presented for payment\\
(3) An amount of Rs. 1,200 for interest on overdraft was debited in the pass book but was intimated to Mr. Zubair on 2\textsuperscript{nd}\\
 July 2022\\
(4) Collection charges of Rs. 28 were entered in cash an amount of Rs. 4,000 but it was wrongly credited to Mr. Yousaf's account by bank\\
(5) On 30\textsuperscript{th}\\
 June 2022 the bank received interest of Rs. 3,000 from a company on Mr. Zubair's deposits with that company but not recorded in cash book.\par
}
\lqn{77}{A trader prepares at the end of each month, a Bank Reconciliation Statement for his business Draft one such statement from under mentioned items as on 31-12-2021 with the help of revised cash book.\\
(1) On 26\textsuperscript{th}\\
 December cheques were issued in favor of creditors for Rs. 38,000 but out of them two cheques for Rs. 28,000 were encashed on 3\textsuperscript{rd}\\
 January 2022\\
(2) One customer had deposited directly into the bank Rs. 5,000 but it was not mentioned in the Cash Book\\
(3) The trader had withdrawn from the bank Rs. 14,000 but the same had not been entered in the Cash Book\\
(4) On 29\textsuperscript{th}\\
 December a cheque for Rs. 12,000 was received and entered in the Cash Bank but had been omitted to be sent to the bank\\
(5) On 27\textsuperscript{th}\\
 December 2021 Rs. 700 were credited in the Pass Book as interest on deposits but the same had been recorded in the Cash Book on 31\textsuperscript{st}  December 2021\\
(6) Bank balance as per the Cash Book Rs. 80,000\par
}
\lqn{78}{From the following particulars prepare a Bank reconcililation Statement of Mr. Irfan as on 31\textsuperscript{st}  March 2022 with the help of a revised cash book:\\
(a) Bank overdraft as per the Cash Book Rs. 24,000\\
(b) Rs. 80,000 worth of cheques were sent for collection and only Rs. 50,000 worth of cheques were collected up to 31\textsuperscript{st}\\
 March\\
(c) Out of cheques for Rs. 60,000 issued in March. Cheques of Rs. 36,000 were presented for payment\\
(d) Bills discounted with the bank and dishonoured Rs. 60,000\\
(e) A cheque for Rs. 10,000 entered in Cash Book (debit side) was not sent to the Bank for collection\\
(f) Discount received Rs. 500 was entered in the Bank Column erroneously.\par
}
\lqn{79}{Prepare a Bank Conciliation Statement under double balance method from the following particulars as on 30-06-2022:\\
(1) Bank balance as per the Bank Statement (Dr.) Rs. 6,000\\
(2) Overdraft as per Cash Book Rs. 7,400\\
(3) Cheques deposited into the bank but not yet collected Rs. 1,940\\
(4) Cheques issued but not yet presented Rs. 3,200\\
(5) The banker charged Rs. 60 for commission\\
(6) Dividend collected by Bank Rs. 200\par
}
\lqn{80}{From the under mentioned particulars, prepare a Bank Reconciliation Statement under double balance method as on 31\textsuperscript{st} March, 2022.\\
(a) Bank Overdraft as per the Bank Statement Rs. 9,200\\
(b) Overdraft as per cash book Rs. 21,200\\
(c) Cheques for Rs. 20,800 were issued in March, out of which those for Rs. 8,000 only were presented up to 31\textsuperscript{st}\\
 March\\
(d) A cheque for Rs. 2,000 was deposited into the bank but was wrongly recorded on the credit side of the Cash Book\\
(e) The credit side of the Cash Book was under cast by Rs. 3,200\\
(f) A cheque for Rs. 1,800 was debited twice in the Cash Book\par
}
\lqn{81}{From the following balances, prepare a trading account for the year ended 31\textsuperscript{st} December 2022. The value of stock on 31\textsuperscript{st} December 2022 was Rs. 34,000.\par
\par\smallskip\noindent\hspace*{1.2em}\begin{tabular}{|cccc|}
\hline
Particulars & Rs. & Particulars & Rs. \\
Purchases & 61,250 & Purchases return & 1,250 \\
Opening Stock & 24,000 & Excise duty & 1,400 \\
Sales & 100,000 & Carriage & 1,250 \\
Sales returns & 1,500 & Clearing charges & 1,600 \\
Octori duty & 1,000 & Wages & 6,000 \\
Gas, water fuel charges & 1,500 & Packing charges & 1,000 \\
Salaries & 7,000 & Custom duty & 1,400 \\
\hline
\end{tabular}\par\smallskip
}
\lqn{82}{(a) Prepare a Trading Account for the year ended 31-12-2022.\\
(b) Prepare a revised Trial Balance after preparation of Trading Account. The following Trial Balance was extracted from the books of Salman on 31-12-2022.\par
\par\smallskip\noindent\hspace*{1.2em}\begin{tabular}{|cccr|}
\hline
Particulars & Dr. & Particulars & C \\
 & Rs. &  & R \\
Opening stock & 116,000 & Capital & 200 \\
Building & 200,000 & Sales & 842 \\
Purchases & 572,000 & Sundry Creditors & 240 \\
Sundry debtors & 278,000 & Purchases Returns & 42 \\
Sales returns & 72,000 &  &  \\
Carriage inward & 12,000 &  &  \\
Carriage outward & 16,000 &  &  \\
Salaries & 32,000 &  &  \\
Furniture & 16,000 &  &  \\
Cash in hand & 10,000 &  &  \\
 & 13,24,000 &  & 13,2 \\
\hline
\end{tabular}\par\smallskip
}
\lqn{83}{Prepare the profit and loss account of Mr. Ajmal for the year ended 30-06-2022 from the following information:\par
\par\smallskip\noindent\hspace*{1.2em}\begin{tabular}{|cccc|}
\hline
Particulars & Rs. & Particulars & Rs. \\
Gross profit & 44,900 & Office electricity charges & 550 \\
Salaries \& Wages & 1,300 & Bad Debts & 640 \\
Insurance & 950 & Office rent & 1,250 \\
Audit fee & 740 & Discount received & 740 \\
Discount allowed & 450 & Dividend received & 900 \\
Bank charges & 1,640 &  &  \\
Commission paid & 400 &  &  \\
\hline
\end{tabular}\par\smallskip
}
\lqn{84}{From the information given below, prepare the profit and loss account of Sadiq \& Co. for the year-ended 31-12-2022.\par
\par\smallskip\noindent\hspace*{1.2em}\begin{tabular}{|ccc|}
\hline
Particulars & Rs. & Particulars \\
Trading A/c (credit balance) & 27,158 & General expenses \\
Salaries to staff & 9,580 & Bad Debts \\
Carriage outward & 300 & Depreciation \\
Rent, Rates \& Taxes & 1,400 & Interset received \\
Selling commission & 268 &  \\
\hline
\end{tabular}\par\smallskip
}
\lqn{85}{From the following Trial Balance prepare a profit and loss account for the year ended 31\textsuperscript{st} December 2022. Also prepare a revised Trial Balance after preparation of profit and loss account.\par
\par\smallskip\noindent\hspace*{1.2em}\begin{tabular}{|cll|}
\hline
 & Dr. & Cr. \\
 & Rs. & Rs. \\
Capital &  & 75,000 \\
Building & 60,000 &  \\
Furniture & 12,000 &  \\
Investment in shares & 16,000 &  \\
Sundry Debtors (A/R) & 24,000 &  \\
Sundry Creditors (A/P) &  & 36,000 \\
Closing Stock & 4,000 &  \\
Gross profit (Credit balance) &  & 32,000 \\
Salaries & 10,000 &  \\
Insurance & 2,400 &  \\
Printing charges & 7,000 &  \\
Advertisement & 8,000 &  \\
Discount allowed & 1,200 &  \\
Discount received &  & 1,000 \\
Bad Debts & 500 &  \\
Interest Received &  & 900 \\
 & 145,100 & 145,100 \\
\hline
\end{tabular}\par\smallskip
}
\lqn{86}{Prepare a trading and profit and loss account for the year ended\\
31\textsuperscript{st} December 2022 and a balance sheet as at that date of Rehman: Closing stock was valued at Rs. 22,500.\par
\par\smallskip\noindent\hspace*{1.2em}\begin{tabular}{|cccrc|}
\hline
 & Debt &  & Credit &  \\
 & Rs. &  & Rs. &  \\
Cash &  & 12,750 &  &  \\
Debtors &  & 26,850 &  &  \\
Freight &  & 675 &  &  \\
Investment &  & 22,500 &  &  \\
Purchases &  & 120,000 &  &  \\
Establishment charges &  & 7,500 &  &  \\
Rent &  & 15,000 &  &  \\
Insurance &  & 4,500 &  &  \\
Delivery Expenses &  & 375 &  &  \\
Wages &  & 7,500 &  &  \\
Sales &  &  &  & 188 \\
Interest on Investment &  &  &  & 90 \\
Creditors &  &  &  & 12,0 \\
Capital &  &  &  & 16,6 \\
 &  & 217,650 &  & 217, \\
\hline
\end{tabular}\par\smallskip
}
\lqn{87}{Prepare a trading and profit and loss account and balance sheet at that date. Closing Stock was valued at Rs. 5,000. Following is the trial balance of Irfan Brothers as on 31\textsuperscript{st}  December 2022:\par
\par\smallskip\noindent\hspace*{1.2em}\begin{tabular}{|cccc|}
\hline
Debt & Rs. & Credit & Rs. \\
Drawings & 3,750 & Bank Overdraft & 12,500 \\
Cash & 22,500 & Sales & 65,500 \\
Goodwill & 6,500 & Capital & 46,250 \\
Purchases & 28,500 & Purchases Return & 750 \\
Stock (1-1-15) & 24,250 & Bills Payable & 7,250 \\
Building & 24,000 &  &  \\
Freight & 250 &  &  \\
Debtors & 15,000 &  &  \\
Trade Expenses & 7,000 &  &  \\
Sales Return & 500 &  &  \\
 & 132,250 &  & 132,250 \\
\hline
\end{tabular}\par\smallskip
}
\lqn{88}{Prepare a trading and profit and loss account for the year ended 30\textsuperscript{th}  June 2022: Closing stock was valued at Rs. 10,000.\par
\par\smallskip\noindent\hspace*{1.2em}\begin{tabular}{|cccrc|}
\hline
 & Debit &  & Credit &  \\
 & Rs. &  & Rs. &  \\
Cash &  & 12,982 &  &  \\
Debtors &  & 3,470 &  &  \\
Creditors &  &  &  & 19,500 \\
Plant &  & 61,300 &  &  \\
Custom duty &  & 8,590 &  &  \\
Salaries &  & 7,850 &  &  \\
Trucks &  & 46,600 &  &  \\
Bank Loan &  &  &  & 16,830 \\
Furniture &  & 36,000 &  &  \\
Sales &  &  &  & 179,900 \\
Purchases &  & 79,200 &  &  \\
Capital &  &  &  & 60,852 \\
Repairs &  & 4,390 &  &  \\
Insurance &  & 16,700 &  &  \\
 &  & 277,082 &  & 277,082 \\
\hline
\end{tabular}\par\smallskip
}
\lqn{89}{Prepare a trading and profit and loss account for the year ended 30-09-2022 and a balance sheet as at that date of Imran. Stock on September 30, 2022 was Rs. 65,000.\par
\par\smallskip\noindent\hspace*{1.2em}\begin{tabular}{|cccrc|}
\hline
 & Debit &  & Credit &  \\
 &  & Rs. & Rs. &  \\
Building &  & 60,500 &  &  \\
Stationary &  & 2,500 &  &  \\
Creditors &  &  &  & 32,500 \\
Commission &  &  &  & 3,750 \\
Taxes &  & 2,500 &  &  \\
Machinery &  & 55,000 &  &  \\
Postage \& Telegram &  & 2,000 &  &  \\
Debtors &  & 62,500 &  &  \\
Drawings &  & 20,000 &  &  \\
Capital &  &  &  & 200,000 \\
Purchases &  & 130,000 &  &  \\
Sales &  &  &  & 165,000 \\
Stock &  & 37,500 &  &  \\
Cash &  & 12,500 &  &  \\
Salaries &  & 8,750 &  &  \\
Carriage &  & 7,500 &  &  \\
 &  & 401,250 &  & 401,250 \\
\hline
\end{tabular}\par\smallskip
}
\lqn{90}{Prepare a trading and profit and loss account for the year ended 30\textsuperscript{th} June 2022 and a balance sheet as at that date of Masood. Closing stock was valued at Rs. 1,250.\par
\par\smallskip\noindent\hspace*{1.2em}\begin{tabular}{|ccccrc|}
\hline
 & Debit &  &  & Credit &  \\
 &  & Rs. &  & Rs. &  \\
Sales &  &  &  &  & 37,50 \\
Sales return &  &  & 1,540 &  &  \\
Purchases &  &  & 32,250 &  &  \\
Purchases return &  &  &  &  & 1,75 \\
Fuel and Power &  &  & 250 &  &  \\
Salaries &  &  & 3,750 &  &  \\
Cash &  &  & 115,875 &  &  \\
Debtors &  &  & 1,875 &  &  \\
Furniture &  &  & 13,125 &  &  \\
Creditors &  &  &  &  & 375 \\
Capital &  &  &  &  & 132,5 \\
Rent &  &  & 1,500 &  &  \\
Commission &  &  & 1,960 &  &  \\
 &  &  & 172,125 &  & 172,1 \\
\hline
\end{tabular}\par\smallskip
}
\lqn{91}{The following Trial Balance was extracted from the books of Saleem on 31-12-2022. Prepare a Trading and profit and loss account for the year ended 31\textsuperscript{st}  December, 2022 and a balance sheet as at that date. Closing Stock was valued at Rs. 3,000.\par
\par\smallskip\noindent\hspace*{1.2em}\begin{tabular}{|cccc|}
\hline
 & Rs. &  & Rs. \\
\underline{Debit Balances:} &  & Sundry Debtors & 56,000 \\
Opening stock & 4,000 & Machinery & 22,000 \\
Purchases & 83,330 & Furniture & 6,000 \\
Carriage inward & 2,680 & Building & 56,000 \\
Wages & 6,900 & Drawings & 6,000 \\
Sales returns & 400 & \underline{Credit Balances:} &  \\
Carriage outward & 790 & Capital & 90,000 \\
General expenses & 640 & Bills Payable & 2,000 \\
Printing charges & 440 & Sundry Creditors & 51,260 \\
Interest Paid & 380 & Purchases Returns & 2,000 \\
Cash in hand & 30,000 & Loan & 16,000 \\
Cash in bank & 36,000 & Discount Received & 300 \\
 &  & Sales & 150,000 \\
\hline
\end{tabular}\par\smallskip
}
\lqn{92}{From the following Trial Balance extracted from the books of Usman \& Sons. Prepare a Trading and profit \& loss account for the year ended 31-12-2022 and a balance sheet as that date. Closing stock was valued at Rs. 40,000.\par
\par\smallskip\noindent\hspace*{1.2em}\begin{tabular}{|cccc|}
\hline
Debit balances & Rs. & Credit balance &  \\
Stock on 01-01-2022 & 60,400 & Sales & 2 \\
Purchases & 231,600 & Capital & 1 \\
Sundry Debtors & 109,660 & Bank Overdraft & 1 \\
Returns Inward & 3,640 & Sundry Creditors & 1 \\
General Expenses & 6,980 &  &  \\
Plant & 22,620 &  &  \\
Wages \& Salaries & 16,740 &  &  \\
Building & 50,000 &  &  \\
Cash in hand & 680 &  &  \\
Cash in bank & 8,720 &  &  \\
Drawings & 16,960 &  &  \\
Motive power & 2,300 &  &  \\
Dock \& Clearing charges & 1,300 &  &  \\
Coal, Gas, Water & 1,700 &  &  \\
Salaries & 9,820 &  &  \\
Interest on overdraft & 4,440 &  &  \\
 & 547,560 &  & 5 \\
\hline
\end{tabular}\par\smallskip
}
\lqn{93}{The balance of provision for Doubtful Debts A/c on 1.1.2022 was Rs. 2200. The bad debts during the year amounted to Rs. 1800. On 31.12.2022 the total balance of Debtors Account was Rs. 37000. The provision for doubtful debts is to be raised to 6\%.\\
Show the journal entries, ledger accounts and profit and loss A/c.\par
}
\lqn{94}{The following particulars are available from the books of Mr. A;\\
Give necessary journal entries and ledger accounts. Also show how items will appear in Balance Sheet.\par
\par\smallskip\noindent\hspace*{1.2em}\begin{tabular}{|c|}
\hline
2022 \\
January 1: Provision for Doubtful Debts \\
Dec. 31: Bad debts \\
Provision for Doubtful Debts to be made at 5\% Debtors on 31st Dec. \\
\hline
\end{tabular}\par\smallskip
}
\lqn{95}{The credit balance of Provision for Doubtful Debts A/c on 1.1.2021 was Rs. 927. Total bad debts during the year were Rs. 1012. At the end of that year, Debtors were amounted to Rs. 25700 and a provision of 5\% was created for doubtful debts. In 2022 actual bad debts amounted to Rs. 800 only. At the end of year, 2022 it was decided to create a provision for doubtful debts for Rs. 100 only.\\
You are required to show the Provision for Doubtful Debts A/c for the above two years.\par
}
\lqn{96}{From the following particulars prepare the relevant accounts including Balance Sheet for the years, 2020, 2021, 2022:\\
Further bad debts Rs. 100 are to be written off in 2015. A provision for doubtful debts is to be raised @ 2\% of sundry debtors:\par
\par\smallskip\noindent\hspace*{1.2em}\begin{tabular}{|cccc|}
\hline
Particulars &  & 2020 & 2021 \\
 &  & Rs. & Rs. \\
Bad debts &  & 400 & 800 \\
Provision for Bad debts on Jan. 1: &  & 2000 & -{}-{}- \\
Bad debts recovered: &  & -{}-{}- & 600 \\
Sundry Debtors on 31\textsuperscript{st} & December: & 70,000 & 30,000 \\
\hline
\end{tabular}\par\smallskip
}
\lqn{97}{The following trial balance was extracted from the books of Rahim on 31-12-2022:\\
Adjustmenrs:\\
(1) Closing stock Rs. 81,000.\\
(2) Building is to be depreciated at the rate of 5\% p.a. and furniture at the rate of 10\% p.a.\\
Prepare trading and profit and loss account and balance sheet.\par
\par\smallskip\noindent\hspace*{1.2em}\begin{tabular}{|lcll|}
\hline
Particulars & Debit & Particulars & Credit \\
 & Rs. &  & Rs. \\
Opening stock & 58,000 & Capital & 100,000 \\
Building & 100,000 & Sales & 421,000 \\
Purchases & 286,000 & Sundry creditors & 120,000 \\
Sundry debtors & 139,000 & Purchases returns & 21,000 \\
Sales returns & 36,000 &  &  \\
Carriage inward & 6,000 &  &  \\
Carriage outward & 8,000 &  &  \\
Salaries & 16,000 &  &  \\
Furniture & 8,000 &  &  \\
Cash in hand & 5,000 &  &  \\
 & 662,000 &  & 662,000 \\
\hline
\end{tabular}\par\smallskip
}
\lqn{98}{From the following trial balance, prepare trading and profit and loss account and balance sheet for the year ended 31-12-2022: Adjustments:\\
(i) Outstanding wages Rs. 1000 and outstanding salaries 1,000.\\
(ii) Depreciation on machinery @ 10\% p.a.\\
(iii) Closing stock 7,000.\par
\par\smallskip\noindent\hspace*{1.2em}\begin{tabular}{|ccc|}
\hline
Particulars & Debit & Credit \\
 & Rs. & Rs. \\
Opening stock & 60,000 &  \\
Factory rent & 1,600 &  \\
Purchases & 21,000 &  \\
Drawings & 7,800 &  \\
Cash & 11,000 &  \\
Machinery & 18,000 &  \\
Salaries & 3,600 &  \\
Debtors & 27,000 &  \\
Wages & 13,000 &  \\
Capital &  & 100,00 \\
Discount & 2,000 & 5,000 \\
Sales &  & 50,000 \\
Creditors &  & 10,000 \\
 & 165,000 & 165,00 \\
\hline
\end{tabular}\par\smallskip
}
\lqn{99}{From the following balances taken from the books of Yaqoobm prepare trading and profit and loss account for the year ended 31\textsuperscript{st} December, 2022 and a balance sheet as on that data:\\
Adjustments:\\
(i) Prepaid rent Rs. 250.\\
(ii) Outstanding commission Rs. 1,250.\\
(iii) Closing stock Rs. 12,500.\\
(iv) Depreciation on furniture @ 7.5\% p.a.\par
\par\smallskip\noindent\hspace*{1.2em}\begin{tabular}{|cccc|}
\hline
Particulars & Rs. & Particulars & Rs. \\
Capital & 27,000 & Rent & 1,250 \\
Returns inward & 1,500 & Furniture & 2,500 \\
Sales & 87,500 & Rates \& taxes & 750 \\
Purchases & 50,000 & Discount allowed & 1,250 \\
Returns outward & 500 & General expenses & 7,500 \\
Commission & 18,750 & Debtors & 14,000 \\
Custom duty & 3,750 & Stock & 11,250 \\
Creditors & 13,250 & Cash in bank & 14,000 \\
 &  & Cash in hand & 1,750 \\
\hline
\end{tabular}\par\smallskip
}
\lqn{100}{From the following trial balance of Arshad, prepare a trading and profit and loss account and balance sheet for the year ended on 31\textsuperscript{st} December 2022:\\
Adjustments:\\
(1) Closing stock on December 31, 2022 was Rs. 105,000.\\
(2) Depreciation on furniture @ 10\% p.a.\\
(3) Rent unexpired Rs. 6,000.\\
(4) Commission earned but not received Rs. 3,000.\par
\par\smallskip\noindent\hspace*{1.2em}\begin{tabular}{|cc|}
\hline
Account Name & Dr. Rs. \\
Cash in Hand & 15,000 \\
Debtors & 42,500 \\
Furniture and Fixtures & 90,000 \\
Opening Stock & 60,000 \\
Goodwill & 60,000 \\
Accumulated Depreciation on Furniture \& Fixtures &  \\
Bank Overdraft &  \\
Arshad's Capital &  \\
Mortage Loan &  \\
Sales Less Returns &  \\
Commission Earned &  \\
Purchase Less Returns & 145,000 \\
Rent Expense & 24,000 \\
Misc. Expense & 2,500 \\
 & 439,000 \\
\hline
\end{tabular}\par\smallskip
}
\lqn{101}{From the following particulars prepare a trading and profit and loss a/c for the year ended on 31\textsuperscript{st} December 2022 and balance sheet as at that date: Adjustments:\\
(1) Closing stock Rs. 17,000.\\
(2) Depreciation motor vehicle 10\%.\\
(3) Outstanding salary Rs. 300.\\
(4) Goods drawn by proprietor for personal use Rs. 600.\par
\par\smallskip\noindent\hspace*{1.2em}\begin{tabular}{|ccc|}
\hline
Particulars & Dr. & Cr. \\
 & Rs. & Rs. \\
Capital &  & 35,000 \\
Motor vehicle & 7,600 &  \\
Drawings & 3,000 &  \\
Purchases & 104,00 &  \\
Freight & 4,300 &  \\
Carriage in & 3,700 &  \\
Returns & 3,200 &  \\
Sales &  & 134,400 \\
Salaries & 6,500 &  \\
Rent & 2,000 &  \\
Stock (1-1-15) & 10,000 &  \\
B/R and B/P & 9,100 & 11,400 \\
Cash in hand & 27,400 &  \\
 & 180,800 & 180,800 \\
\hline
\end{tabular}\par\smallskip
}
\lqn{102}{From the following trial balance of Mr. Arif. Prepare a trading and profit and loss account and a balance sheet for the year ended 31\textsuperscript{st} December 2022.\\
Adjustments:\\
(1) Stock on 31\textsuperscript{st}  December 2022 was Rs. 2,900\\
(2) Depreciation plant and machinery by Rs. 800.\\
(3) repairs Rs. 80 is outstanding.\\
(4) Make a provision of 10\% on debtors for bad and doubtful debts.\par
\par\smallskip\noindent\hspace*{1.2em}\begin{tabular}{|ccl|}
\hline
Particulars & Rs. & R \\
Capital &  & 20 \\
Creditors &  & 3 \\
Returns & 600 & 1 \\
Sales &  & 32 \\
Plant \& Machinery & 8,000 &  \\
Sundry debtors & 4,800 &  \\
Drawings & 2,000 &  \\
Purchases & 21,000 &  \\
Bank & 10,000 &  \\
Repairs & 100 &  \\
Stock (01-01-15) & 4,000 &  \\
Rent & 800 &  \\
Manufacturing Expenses & 3,600 &  \\
Trade expenses & 1,400 &  \\
Bad debts & 400 &  \\
Fuel \& Power & 500 &  \\
 & 57,200 & 57 \\
\hline
\end{tabular}\par\smallskip
}
\lqn{103}{Following is the Trial Balance of Mukhtar Trading Co. as on\\
31\textsuperscript{st} December 2015. You are required to prepare the Trading and profit and loss account of the business for the year ending 31\textsuperscript{st} December 2022 and a balance sheet as at that date.\\
Adjustments:\\
(a) Wages payable Rs. 500.\\
(b) Bad debts provision to be adjusted to 2\% on Debtors.\\
(c) Allow interest on capital at 5\%.\\
(d) Closing stock valued at Rs. 15,000.\par
\par\smallskip\noindent\hspace*{1.2em}\begin{tabular}{|ccc|}
\hline
Particulars & Rs. & Rs. \\
Opening stock & 10,000 &  \\
Purchases and sales & 40,000 & 100,000 \\
Sundry Debtors \& Creditors & 30,000 & 8,125 \\
Plant \& Machinery & 35,000 &  \\
Freehold Premises & 24,000 &  \\
Salaries & 5,000 &  \\
Wages & 6,575 &  \\
Carriage Outward & 375 &  \\
Printing \& Stationary & 400 &  \\
Bad Debts & 500 &  \\
Provision for Bad debts &  &  \\
Trade Expenses & 750 & 225 \\
Cash in Hand & 5,250 &  \\
Capital &  & 50,000 \\
Insurance & 500 &  \\
 & 158,350 & 158,350 \\
\hline
\end{tabular}\par\smallskip
}
\lqn{104}{The following Trial Balance has been extracted from the books of a Trader on 31\textsuperscript{st} March, 2022. Prepare Trading and profit and loss account for the year ending 31\textsuperscript{st} March, 2022 and a Balance Sheet as at that date. Adjustments:\\
(a) Reserve for discount 2.5\% both on debtors and creditors.\\
(b) Depreciate horses and carts at 7.5\%.\\
(c) Stock on 31\textsuperscript{st}  March 2022 Rs. 120,000.\\
(d) Interest accrued Rs. 500.\par
\par\smallskip\noindent\hspace*{1.2em}\begin{tabular}{|ccc|}
\hline
Particulars & Dr. & Cr \\
 & Rs. & Rs \\
Opening stock & 265,000 &  \\
Debtors & 65,000 &  \\
Horses and carts & 30,000 &  \\
Purchases & 300,000 &  \\
General Expenses & 26,000 &  \\
Traveling Expenses & 21,100 &  \\
Commission & 19,500 &  \\
Cahs in hand & 60,500 &  \\
Printing charges & 13,400 &  \\
Interest on Mortgage & 2,000 &  \\
Building & 91,000 &  \\
Capital &  & 211, \\
Sales &  & 577,4 \\
Mortgage loan &  & 62,0 \\
Sundry creditors &  & 42,0 \\
Interest &  & 1,0 \\
 & 893,500 & 893,5 \\
\hline
\end{tabular}\par\smallskip
}
\lqn{105}{From the following Trial Balance of Azan \& Co. You are required to prepare Trading and profit and loss account for the year ended\\
31\textsuperscript{st} December, 2022 and a Balance sheet as at that date: Adjustments:\\
(a) Stock on 31-12-2022 was valued at Rs. 40,000.\\
(b) Depreciation Plant \& Machinery at 10\% p.a.\\
(c) Provide 5\% on debtors for Bad debts and 2\% for discount.\par
\par\smallskip\noindent\hspace*{1.2em}\begin{tabular}{|ccc|}
\hline
Particulars & Dr. & Cr. \\
 & Rs. & Rs. \\
Opening Stock & 60,000 &  \\
Goodwill & 70,000 &  \\
Plant \& Machinery & 105,000 &  \\
Loose Tools & 1,000 &  \\
Capital &  & 120,000 \\
Bills Receivable \& Bills Payable & 5,000 & 16,000 \\
Purchases and Sales & 90,000 & 228,500 \\
Rent & 1,000 &  \\
Advertisement & 28,000 &  \\
Sundry Debtors \& Creditors & 40,000 & 60,000 \\
Cash in hand & 21,500 &  \\
Furniture & 2,000 &  \\
Returns & 2,500 & 3,000 \\
Bad Debts & 1,500 &  \\
 & 427,500 & 427,500 \\
\hline
\end{tabular}\par\smallskip
}
\lqn{106}{From the following Trial Balance of Shazaib \& Co. You are required to prepare Trading and Profit and loss account for the year ending\\
31\textsuperscript{st} December, 20225 and a Balance Sheet as at that:\\
Adjustments:\\
(a) Closing stock was valued at Rs. 100,000.\\
(b) Depreciate, Land \& Building at 2.5\%.\\
(c) Write off Rs. 4,800 bad debts from debtors.\\
(d) Maintain Reserve for doubtful debts at 8\%.\par
\par\smallskip\noindent\hspace*{1.2em}\begin{tabular}{|ccc|}
\hline
Particulars & Debit & Cr \\
 & Rs. & R \\
Land \& Building & 250,000 &  \\
Furniture \& Fixtures & 26,000 &  \\
Opening Stock & 80,600 &  \\
Drawings and Capital & 64,000 & 280 \\
Purchases \& Sales & 220,000 & 408 \\
Debtors \& Creditors & 60,800 & 48 \\
Returns & 12,000 & 10 \\
Loan from SBFC &  & 80 \\
Wages \& Salaries & 53,000 &  \\
Printing and Stationary & 6,200 &  \\
Bad debts & 8,400 &  \\
Provision for doubtful debts &  & 6, \\
Cash at Bank & 51,760 &  \\
 & 832,760 & 832 \\
\hline
\end{tabular}\par\smallskip
}
\lqn{107}{The following Trial Balance was extracted from the books of Saad \& Bros on 31\textsuperscript{st} December, 2022. You are requested to prepare trading and profit and loss account for the year ending 31-12-2022 and a Balance Sheet at that date.\\
Adjustments:\\
(a) Provision for bad debts is to be maintained at 5\% on sundry debtors.\\
(b) Closing stock was valued at Rs. 59,250.\\
(c) The manager is entitled to a communication of 10\% on the net profit after charging such commission.\par
\par\smallskip\noindent\hspace*{1.2em}\begin{tabular}{|cccc|}
\hline
Debit Balances & Rs. & Debit Balances & Rs. \\
Drawings & 6,000 & Bad debts & 2,050 \\
Cash in Hand & 550 & Sundry Debtors & 34,000 \\
Cash at Bank & 8,750 & Bills Receivable & 3,975 \\
Furniture & 9,000 & \underline{Credit Balances:} &  \\
Plant \& Machinery & 15,000 & Discount received & 2,550 \\
Opening Stock & 30,000 & Sales (less returns) & 164,000 \\
Purchases (Less returns) & 115,500 & Provision for Bad debts & 1,000 \\
Transportation in & 15,000 & Sundry creditors & 41,250 \\
Rent \& Taxes & 11,275 & Capital & 49,300 \\
Discount allowed & 3,400 &  &  \\
Repairs and Renewals & 3,600 &  &  \\
\hline
\end{tabular}\par\smallskip
}
\lqn{108}{From a following trial balance of Anwer prepare a Trading and profit and loss account for the year ending 31\textsuperscript{st} December 2022 and balance sheet as on that after making neessary adjustments.\\
Adjustments:\\
(a) Closing stock was valued at Rs. 55,000.\\
(b) Depreciation on Plants \& Machinery at 10\%.\\
(c) Prepaid salaries Rs. 2,400.\\
(d) General Manager is entitled a commission of 10\% before charging such commission.\par
\par\smallskip\noindent\hspace*{1.2em}\begin{tabular}{|cccc|}
\hline
Debit Balances & Rs. & Credit Balance &  \\
Opening stock & 34,000 & Sales & 10 \\
Purchases & 71,500 & Capital & 11 \\
Bills receivable & 14,700 & Bills Payable & 1 \\
Sundry debtors & 18,000 & Sundry creditors & 7 \\
Freight & 4,750 & Commission & 8 \\
Octori duty & 550 & Discount & 3 \\
Salaries & 13,050 & Interest & 1 \\
General expenses & 3,550 & Rent & 3 \\
Discount & 2,470 &  &  \\
Rent \& Rates & 5,700 &  &  \\
Custom Duty & 2,900 &  &  \\
Cash in hand & 15,500 &  &  \\
Goodwill & 3,750 &  &  \\
Income Tax & 2,700 &  &  \\
Drawings & 5,000 &  &  \\
Plant \& Machinery & 58,380 &  &  \\
 & 256,500 &  & 25 \\
\hline
\end{tabular}\par\smallskip
}
\lqn{109}{The following balances are extracted from the Books of Asif \& Bros. For the year ended 31\textsuperscript{st}  December 2022. You are required to prepare the work sheet for the year ended 31-12-2022.\\
The Merchandise inventory on 31-12-2022 was valued at Rs. 53,000.\par
\par\smallskip\noindent\hspace*{1.2em}\begin{tabular}{|cc|}
\hline
 & Rs. \\
Merchandise inventory on 01-01-2015 & 36,000 \\
Perchases & 143,600 \\
Purchases discounts & 3,000 \\
Purchases returns and allowances & 2,500 \\
Sales & 175,000 \\
Wages & 1,500 \\
Sales salaries & 15,000 \\
Advertising expenses & 9,600 \\
Utilities expenses & 8,900 \\
Insurance & 950 \\
Rent, rates and taxes & 2,300 \\
\hline
\end{tabular}\par\smallskip
}
\lqn{110}{The following Trial Balance was extracted from the books of Baber \& Sons on 31\textsuperscript{st}  March, 2022. You are required to prepare a worksheet for the year ended on 31-03-2022.\\
Additional Information:\\
(1) Outstanding salaries Rs. 1,000.\\
(2) Depreciation on plant and machinery @ 10\% p.a.\\
(3) Merchandise inventory on 31-03-2015 was valued at Rs. 6,000.\par
\par\smallskip\noindent\hspace*{1.2em}\begin{tabular}{|ccc|}
\hline
Accounts Name & Dr. Rs. & C \\
Cash & 5,000 &  \\
Account Receivable & 16,400 &  \\
Merchandise inventory on 01-04-2014 & 6,000 &  \\
Plant \& Machinery & 106,000 &  \\
Furniture \& Fixtures & 2,600 &  \\
Capital &  & 1 \\
Accounts Payable &  &  \\
Purchases & 60,000 &  \\
Sales &  & 7 \\
Sales returns and allowances & 3,000 &  \\
Sales Discount & 1,600 &  \\
Salaries expenses & 12,000 &  \\
 & 212,600 & 2 \\
\hline
\end{tabular}\par\smallskip
}
\lqn{111}{From the following Trial Balance of Ahmed \& Co. Prepare a ten column work sheet for the year ended 31\textsuperscript{st} December 2022:\\
Adjustments:\\
(1) Merchandise inventory on December 31, 2022 was Rs. 105,000.\\
(2) Depreciation on furniture @ 10\% p.a.\\
(3) Rent unexpired Rs. 6,000.\\
(4) Commission earned but not received Rs. 3,000.\par
\par\smallskip\noindent\hspace*{1.2em}\begin{tabular}{|ccc|}
\hline
Accounts Name & Dr. Rs. & Cr. Rs. \\
Cash in hand & 15,000 &  \\
Account Receivables & 42,500 &  \\
Furniture \& Fixtures & 90,000 &  \\
Merchandise inventory (01-01-2015) & 60,000 &  \\
Goodwill & 60,000 &  \\
Accumulated Depreciation on Furniture \& Fixture &  & 9,000 \\
Bank Overdraft &  & 27,500 \\
Ahmad's Capital &  & 95,000 \\
Mortgage loan &  & 60,000 \\
Sales less returns &  & 232,500 \\
Commission earned &  & 15,000 \\
Purchases less returns & 145,00 &  \\
Rent Expense & 24,000 &  \\
Mics. Expenses & 2,500 &  \\
 & 439,000 & 439,000 \\
\hline
\end{tabular}\par\smallskip
}
\lqn{112}{The following Trial Balance has been extracted from the books of Usman Bros for the year ended on 31\textsuperscript{st} December 2022.\\
Adjustments:\\
(1) Closing inventory was valued at Rs. 14,000.\\
(2) Advertisement prepaid Rs. 1,790.\\
(3) Accrued dividend Rs. 120.\\
(4) Allow interest on capital @ 5\% p.a.\par
\par\smallskip\noindent\hspace*{1.2em}\begin{tabular}{|ccc|}
\hline
Names of Accounts & Dr. Rs. & Cr. \\
Cash in hand & 435 &  \\
Prepaid advertisement & 2,790 &  \\
Office Furniture & 25,400 &  \\
Accumulated Depreciation &  & 1, \\
Capital &  & 30 \\
Drawings & 2,600 &  \\
Accounts receivable & 29,500 &  \\
Accounts Payable &  & 8, \\
Sales &  & 66 \\
Salaries & 9,850 &  \\
Purchases & 34,980 &  \\
Transporting Expenses & 1,745 &  \\
Dividend Received &  & 9 \\
 & 107,300 & 107 \\
\hline
\end{tabular}\par\smallskip
}
\lqn{113}{From the following balances prepare ten columns work sheet for the year ended 31\textsuperscript{st}  December 2022.\\
Adjustments:\\
(1) Closing inventory Rs. 102,000.\\
(2) Write off Rs. 500 as uncollectible and make an allowance at 5\% on accounts receivable for uncollectible.\\
(3) Commission unearned Rs. 1,500.\par
\par\smallskip\noindent\hspace*{1.2em}\begin{tabular}{|ccc|}
\hline
Names of Accounts & Dr. Rs. & Cr. Rs. \\
Accounts Receivable & 52,775 &  \\
Land \& Buildings & 50,000 &  \\
Plant \& Machinery & 42,000 &  \\
Wages Productive & 5,620 &  \\
Opening Inventory & 69,680 &  \\
Travelling Expenses & 9,750 &  \\
Repair \& Renewal & 1,430 &  \\
Uncollectible & 1,240 &  \\
Salaries & 14,800 &  \\
Capital &  & 114,500 \\
Bank Overdraft &  & 56,935 \\
Sales &  & 189,000 \\
Commission &  & 5,500 \\
Accounts Payable &  & 31,000 \\
Purchases & 149,640 &  \\
 & 396,935 & 396,935 \\
\hline
\end{tabular}\par\smallskip
}
\lqn{114}{Ascertain the cost of goods sold from the following particulars:\par
\par\smallskip\noindent\hspace*{1.2em}\begin{tabular}{|ccc|}
\hline
(i) & Opening Merchandise Inventory & 19 \\
(ii) & Purchases (gross) & 15 \\
(iii) & Purchases Discounts & 3 \\
(iv) & Purchases returns and allowances & 2 \\
(v) & Freight Inwards & 2 \\
(vi) & Custom duty and clearing charges & 2 \\
(vii) & Closing Merchandise Inventory & 30 \\
\hline
\end{tabular}\par\smallskip
}
\lqn{115}{Ascertain the gross income from the following particulars.\par
\par\smallskip\noindent\hspace*{1.2em}\begin{tabular}{|cc|}
\hline
 & Rs. \\
Purchases during the year, 2015 & 350,000 \\
Purchases discount & 3,700 \\
Purchases returns and allowances & 2,050 \\
Sales & 470,000 \\
Sales discount & 1,500 \\
Sales returns and allowances & 2,200 \\
Merchandise inventory on 01-01-2015 & 50,500 \\
Merchandise inventory on 31-12-2015 & 30,800 \\
Transportation inwards & 9,500 \\
\hline
\end{tabular}\par\smallskip
}
\lqn{116}{The following balance are extracted from the books of Asif \& Bros. For the year ended 31\textsuperscript{st} December 2022, you are required to prepare an income statement for the year ended 31-12-2022.\\
The Merchandise inventory on 31-12-2022 wsa valued at Rs. 53,000.\par
\par\smallskip\noindent\hspace*{1.2em}\begin{tabular}{|cc|}
\hline
Merchandise inventory on 01-01-2015 & 36 \\
Purchases & 14 \\
Sales & 17 \\
Wages & 1 \\
Sales salaries & 15 \\
Advertising expenses & 9 \\
Utilities expenses & 7 \\
Insurance & 9 \\
Rent, rates and taxes & 2 \\
Commission to salesman & 1 \\
\hline
\end{tabular}\par\smallskip
}
\lqn{117}{The following Trial Balance was extracted from the books of Naeem \& Sons on 31\textsuperscript{st} March 2022. From this you are required to prepare an income statement for the year ended on 31-03-2022 and a Balance Sheet as at that date.\\
Additional Information:\\
(1) Outstanding salaries Rs. 1,000.\\
(2) Depreciation on Plants \& Machinery @ 10\% p.a.\\
(3) Merchandise inventory on 31-03-20225 was valued at Rs. 6,000.\par
\par\smallskip\noindent\hspace*{1.2em}\begin{tabular}{|ccc|}
\hline
Particulars & Dr. Rs. & Cr. Rs. \\
Cash & 5,000 &  \\
Accounts Receivable & 9,000 &  \\
Merchandise inventory on 01-04-2014 & 6,000 &  \\
Plant \& Machinery & 24,000 &  \\
Land \& Building & 94,000 &  \\
Furniture \& Fixtures & 2,600 &  \\
Capital &  & 138,800 \\
Accounts Payable &  & 3,800 \\
Purchases & 60,000 &  \\
Sales &  & 70,000 \\
Salaries expenses & 12,000 &  \\
 & 212,600 & 212,600 \\
\hline
\end{tabular}\par\smallskip
}
\lqn{118}{The following are the balance extracted from the books of Usman Traders as at 31\textsuperscript{st} December 20225.\\
Adjustments:\\
(1) Merchandise inventory at 31-12-2022 was valued at Rs. 25,470.\\
(2) Insurance unexpired Rs. 2,140.\\
(3) Allow interest on capital at 5\% p.a. \& 3\% on drawings.\\
Prepare income statement and balance sheet.\par
\par\smallskip\noindent\hspace*{1.2em}\begin{tabular}{|cccc|}
\hline
Debit Balances & Rs. & Credit Balances &  \\
Inventory (01-01-2022) & 45,270 & Capital & 5 \\
Drawings & 4,550 & Accounts payable & 4 \\
Buildings & 400,000 & Allowances for uncollectible &  \\
Office Equipment & 100,000 & Discount &  \\
Cash in Hand & 29,150 & Sales & 3 \\
Cash in Bank & 116,540 & Purchases returns &  \\
Accounts Receivable & 69,750 &  &  \\
Purchases & 232,710 &  &  \\
Sales Returns & 5,540 &  &  \\
Salaries & 37,875 &  &  \\
Travelers Commission & 3,940 &  &  \\
Office Expenses & 4,500 &  &  \\
Insurance & 4,270 &  &  \\
 & 10,54,095 &  & 10 \\
\hline
\end{tabular}\par\smallskip
}
\lqn{119}{On 31\textsuperscript{st} December 2022, the following Trial Balance was extracted from the books of Khurram traders.\\
Adjustments:\\
(1) Taxes outstanding Rs. 2,000; Insurance unexpired Rs. 500.\\
(2) Reserve for discount on debtors and creditors at 2\%.\\
(3) Goods costing Rs. 5,000 were taken by Khurram for his personal use, for which no record has been kept.\\
Prepare Income Statement and Balance Sheet.\par
\par\smallskip\noindent\hspace*{1.2em}\begin{tabular}{|cccc|}
\hline
Particulars & Rs. & Particulars & Rs. \\
Furniture & 65,000 & Bought ledger balance & 61,750 \\
Sales Ledger Balance & 74,750 & Income tax & 950 \\
Merchandise inventory on 31-12-2015 & 42,500 & Bank overdraft & 18,350 \\
Purchases Adjustment & 181,400 & Sales & 296,450 \\
Cash in hand & 8,220 & Capital & 100,950 \\
Freight & 4,370 & Sales salaries & 24,160 \\
Taxes \& Insurance & 8,000 & Accountancy fee & 14,150 \\
Carriage on sales & 2,910 &  &  \\
Utilities charges & 51,090 &  &  \\
\hline
\end{tabular}\par\smallskip
}
\lqn{120}{Following in the Trial Balance of a firm. From this you are required to prepare an income statement for the year ending on 31\textsuperscript{st} December, 2022 and a Balance Sheet as on that date.\\
Adjustments:\\
1. Merchandise inventory on 31.12.2022 was valued at Rs. 35,000.\\
2. Store supplies on hand 31.12.2022 Rs. 1,500.\\
3. Depreciate furniture and fixture and plant and Machinery by 10\% p.a.\\
4. Write off bad debts from accounts receivable Rs. 1,500.\par
\par\smallskip\noindent\hspace*{1.2em}\begin{tabular}{|ccrr|}
\hline
\multicolumn{4}{|c|}{Trial Balance On 31.12.2015} \\
 &  & Dr. &  \\
 &  & Rs. &  \\
Cash at Bank &  & 95,500 &  \\
Accounts receivable &  & 35,200 &  \\
Stores supplies &  & 5,000 &  \\
Prepaid Rent &  & 11,240 &  \\
Furniture and Fixtures &  & 7,600 &  \\
Accumulated depreciation- Furniture and Fixture &  &  &  \\
Plant and Machinery &  & 45,000 &  \\
Accumulated depreciation- Plant and Machinery &  &  &  \\
Capital &  &  &  \\
Accounts payable &  &  &  \\
Drawings &  & 31,000 &  \\
Sales revenue &  &  & 2 \\
Salaries expense &  & 9,500 &  \\
Purchases &  & 95,000 &  \\
Wages &  & 10,000 &  \\
Purchases returns &  &  &  \\
Sales returns &  & 3,000 &  \\
Merchandise Inventory on 1.1.2015 &  & 45,000 &  \\
Traveling expenses &  & 5,460 &  \\
Miscellaneous expenses: &  & 5,000 &  \\
 &  & 403,500 & 40 \\
\hline
\end{tabular}\par\smallskip
}
\lqn{121}{State with reason whether the following items of expenditures are capital or revenue:\\
(a) Carriage paid on goods purchased.\\
(b) Wages paid to workmen for construction of a new building.\\
(c) Replacement cost of a worn-out part of plant.\\
(d) Repair to furniture purchased second hand.\\
(e) Addition and alternation made in the building.\par
}
\lqn{122}{Show by giving reasons, whether the following items of expenditures are capital or revenue:\\
(i) Expenditure on advertising campaign.\\
(ii) Amount spent on the uniform of the staff.\\
(iii) Preliminary expenses incurred on the flotation of a company.\\
(iv) Freight paid on a newly purchased machine.\\
(v) An amount spent as legal expenses for abuse of trade mark.\par
}
\lqn{123}{State with reasons which item of expenditure would be charged to capital or revenue:\\
(a) Cost of a new machinery.\\
(b) Legal expenses incurred in an income tax appeal.\\
(c) Cost of air conditioning the office of the general manager.\\
(d) Cost of acquiring the goodwill of an old firm.\\
(e) Cost of re-decorating a cinema hall.\par
}
\lqn{124}{State with reasons whether the following expenditures are capital or revenue:\\
(i) Repair to a motor car met with an accident.\\
(ii) Painting expenditures of a newly constructed factory.\\
(iii) Compensation paid to a person injured by company's car.\\
(iv) Expenditure on replacement of state roof by glass roof.\\
(v) Cost of goodwill.\par
}
\lqn{125}{Would you consider the following as revenue or capital charges? State your reasons in each case:\\
(a) Repairs and renewal of machinery.\\
(b) Cost incurred in acquiring a patent.\\
(c) Cost constructing an extra door in the building.\\
(d) Compensation paid to workers in termination of their services.\\
(e) Cost incurred by a firm on advertising a new product.\par
}
\lqn{126}{State with reasons whether the following expenditures should be capital or revenue.\\
(i) Cost of plant and machinery.\\
(ii) A heavy advertising expenditure was made to introduce a new product.\\
(iii) Wages paid to workers for manufacturing goods.\\
(iv) Purchase of lease-hold land.\\
(v) Discount on issue of debentures.\par
}
\lqn{127}{State with reason whether the following items should be regarded as capital or revenue:\\
(i) Wages paid for the installation of machinery.\\
(ii) Payment made on purchase of furniture.\\
(iii) Custom duty paid on machinery.\\
(iv) Cost of air conditioning.\\
(v) Profit earned on sale of fixed assets.\par
}
\lqn{128}{State with reasons whether the following items of expenditures are capital, revenue or deferred revenues:\\
(i) Medicines were purchased for Rs. 50,000 and Rs. 1,000 were spent on carrying them to hospital.\\
(ii) Boundary walls were constructed by a college at a cost of Rs. 30,000.\\
(iii) Freight and cartage on the new machine Rs. 200 and errection charges Rs. 500.\\
(iv) Rs. 5,000 were paid for the removal of stock from old works to new works.\\
(v) Commission on issue of debentures.\par
}
\lqn{129}{Explain with reasons whether the following items relating to a Sugar Mill Company are capital or revenue:\\
(a) Rs. 25,000 were the amount of premium was received on the issue of shares.\\
(b) Labour welfare expenses incurred Rs. 20,000.\\
(c) Rs. 300,000 paid for excise duty on goods manufactured.\\
(d) A machinery standing in the books at Rs. 300,000 were sold for Rs. 240,000.\\
(e) Legal expenses incurred in an action for infringement of its trade mark.\par
}
\lqn{130}{Explain by giving reasons whether the following items of expenditures and receipts are capital or revenue:\\
(i) Profit on sale of shares held by a firm.\\
(ii) Amount spent on the purchase of patent right.\\
(iii) Interest paid on loan borrowed for business.\\
(iv) Premium given for lease.\\
(v) Legal expenses incurred on income tax appeal.\par
}
\lqn{131}{State which items are chargeable to capital and which to revenue:\\
(a) Purchased goods for resale purposes for Rs. 30,000.\\
(b) Old machinery was repaired at a cost of Rs. 2,000.\\
(c) Old machinery having book value of Rs. 40,000 was sold for Rs. 37,000.\\
(d) Amount spent on the uniform of factory workers.\\
(e) Compensation paid to an employee for injury caused during his employment.\par
}
\lqn{132}{Would you consider the following chargeable to Capital or Revenue? Give reasons for your answers.\\
(a) Cost incurred for obtaining a mortgage against a long-term debt.\\
(b) Premium given for lease.\\
(c) Commission on issue of debentures.\\
(d) Discount on issue of debentures.\\
(e) Expenditure incurred on renovation of the old building.\\
(f) Cost of registration of newly-formed limited company.\\
(g) Legal expenses incurred for checking abuse of our trade mark by others.\\
(h) Expenditure incurred on the construction of temporary huts which were necessary for the construction of a cinema hall.\\
(i) Repair necessitated by negligence and lack of care and attention.\\
(j) Cost of experiment which did not prove success.\par
}
\lqn{133}{Classify the following items into capital and revenue:\\
(a) Purchase of furniture\\
(b) Repair to machinery\\
(c) Payment of salaries\\
(d) Whitewashing expenses\\
(e) Depreciation\\
(f) Purchase of newspapers\\
(g) Sale of old furniture\\
(h) Sale of old newspapers\\
(i) Advertising campaign to launch a new product\\
(j) Repair to machinery purchased second-hand\\
(k) Cost of erecting new machinery\\
(l) Extension of building\\
(m) Registration expenses for purchase of land\\
(n) Expenses for removal of factory from old place to a new and better place\par
}
\lqn{134}{Explain by giving reasons whether the following items of expenditure and receipt and Capital or Revenue.\\
(i) Profit on sale of shares held by a firm.\\
(ii) Replacement expenditure of a worn out part of machinery.\\
(iii) Compensation paid to an employee for injury caused during his employment.\\
(iv) Legal expenses incurred in an income tax appeal.\\
(v) Wages paid to workers for manufacturing goods.\\
(vi) Amount spent on the purchase of patent right.\\
(vii) Interest paid on loan borrowed for business.\\
(viii) Amount spent on the uniform of factory workers.\\
(ix) Purchase of lease-hold land.\\
(x) Preliminary expenses incurred on the formation of a new company.\\
(xi) Legal expenses in purchasing landing property.\\
(xii) Regular fee paid to the legal Advisor of the company.\par
}
\lqn{135}{How would you rectify the following errors detected before preparing the Trial Balance:\\
(a) Sales returns book has been under-cast by Rs. 1,000.\\
(b) A payment of Rs. 8,000 on account of rent has been posted twice to Rent A/c.\\
(c) A sum of Rs. 4,000 received from Nasir has not been posted to his account.\\
(d) A sum of Rs. 1,500 debited to Imran A/c as Rs. 150.\\
(e) A sale of Rs. 6,000 has been credited to Customer’s A.c.\par
}
\lqn{136}{Rectify the following errors:\\
(a) Sales book under-cast by Rs. 10,000.\\
(b) Purchases book under-cast by Rs. 5,000.\\
(c) Sales book overcast by Rs. 2,000.\\
(d) Purchases book overcast by Rs. 3,000.\par
}
\lqn{137}{Rectify the following errors, if they are detected after the preparation of the Trial balance:\\
(a) Return inward book has been overcast by Rs. 400.\\
(b) Return outward book has been overcast by Rs. 900.\\
(c) Return inward has been under-cast by Rs. 3,200.\\
(d) Return outward has been under-cast by Rs. 2,600.\\
(e) Rent paid Rs. 4,000 debited to land lord’s personal account.\par
}
\lqn{138}{Rectify the following errors:\\
(i) Sales to A Rs. 5,000 credited to his account.\\
(ii) Sales to A Rs. 5,000 debited to his account as Rs. 5000.\\
(iii) Purchases from B Rs. 4,600 credited to his account as Rs. 460.\\
(iv) Purchases from B Rs. 4,600 debited to his account.\\
(v) Purchases from B Rs. 2,000 credited to his account as Rs. 3,000.\par
}
\lqn{139}{Give entries to correct the errors:\\
(i) A bill for Rs. 2,500 for furniture sold to B. Brother’s was credited to sales account.\\
(ii) An amount of Rs. 1,780 received on account of interest was credited to commission account.\\
(iii) Goods to the value of Rs. 2,600 were returned by Ali, but return was not recorded in the books.\\
(iv) Wages paid Rs. 2,000 posted to the wages a/c twice.\\
(v) The amount of Rs. 500 written off bad debts in the previous year now received credited to the personal a/c of the debtors.\par
}
\lqn{140}{Give journal entries to correct the errors:\\
(i) An item of Rs. 3,000 paid for the purchase of furniture debited to purchases account.\\
(ii) An amount Rs. 600 withdrawn by the proprietor for his personal use has been debited to trade expenses a/c.\\
(iii) Salary Rs. 1,000 paid to clerk debited to his personal a/c.\\
(iv) A bill for the extension of premises amounting to Rs. 5,100 has been debited to repair account.\\
(v) Purchases from X Rs. 4,600 debited to his account.\par
}
\lqn{141}{Pass entries for the rectification of errors given below:\\
(i) Legal expenses Rs. 6,000 paid to lawyer have been wrongly debited to his personal account.\\
(ii) Wages paid Rs. 6,000 wrongly debited to custom duty.\\
(iii) Trade expenses of Rs. 360 posted in the ledger as Rs. 810.\\
(iv) A sale of Rs. 6,000 to A was wrongly debited to the account of B.\\
(v) Sales book was overcast by Rs. 6,000.\par
}
\lqn{142}{Rectify the following errors by passing necessary journal entries:\\
(i) A sale of Rs. 992 to Ashraf was entered in the sale ledger as Rs. 1,388. Ashraf A/c was however correctly debited.\\
(ii) A bill received from Salman \& Co. for repairs done to furniture Rs. 400 was wrongly entered in the Purchases Book.\\
(iii) An item of Rs. 8,000 relating to prepaid rent was omitted to be brought forward.\\
(iv) An item of purchase of Rs. 3,0200 was entered in the Purchases Account at Rs. 300.\\
(v) A sale return of Rs. 2,000 were not passed through the sale returns book, though goods were duly taken into the stock.\par
}
\lqn{143}{The following errors were found from the Books of a firm. Pass the necessary entries to rectify them:\\
(a) Rs. 168,000 paid for Wages to workmen for making show cases had been charged to the Wages A/c.\\
(b) A sale of Rs. 36,000 has been passed through the Purchases Book. The customer’s account has, however been correctly debited.\\
(c) A sale of Rs. 9,000 to Kabir Bros. was credited to their account.\\
(d) Rs. 60,000 paid for the purchase of a Motor Cycle for a partner, had been charged to the miscellaneous expenses account.\\
(e) A purchase of Rs. 1,392 had been debited to Supplier’s account as Rs. 1,264. The supplier was Sarwar.\par
}
\lqn{144}{Rectify the following errors:\\
(i) Rs. 1,680, as debit balance in Amjad account are written off as bad debts recorded in Amjad account only.\\
(ii) Rs. 4,000 paid for college fees of proprietor’s son were charged to the Trade expenses account.\\
(iii) Bad debts recovered Rs. 1,000 from Arshad were credited to their account.\\
(iv) A receipt of Rs. 3,500 from Saeed \& Co. was credited to their account as Rs. 5,300.\\
(v) Sales returns of Rs. 1,000 were taken into Stock but no entry was passed in the books.\par
}
\lqn{145}{Pass journal entries to rectify them and prepare the suspense account:\\
(i) The total of the credit side of Mohsin A/c was overcast by Rs. 2,000.\\
(ii) Sales Book was under-cast by Rs. 2,000.\\
(iii) The total of purchases return book Rs. 4,000 was not posted to the legder.\\
(iv) A credit balance of Rs. 7,550 of the rent received account was shown as Rs. 5,700.\\
(v) Goods amounting to Rs. 1,500 were sold to Mahmood \& Co. for cash. It was correctly recorded in the Cash Book, but was wrongly posted to mahmood \& Co. account also.\\
(vi) The discount column on debit side of cash book has been posted to the credit side of discount account Rs. 80.\par
}
\lqn{146}{Give journal entries to rectify the following errors:\\
(i) A purchase of Rs. 20,000 from Ali was wrongly passed through the Sale Book.\\
(ii) A purchase of Rs. 14,000 from Danish was wrongly passed through the purchases Returns Book.\\
(iii) A purchase of Rs. 9,600 from Saad was wrongly passed through the Sales Returns Book.\\
(iv) A sale of Rs. 16,000 to Babar was wrongly passed through the Sales Returns Book.\\
(v) A sale of Rs. 18400 to Liaqat was wrongly passed through the Purchases Returns Book.\\
(vi) A cheque received from Ali of Rs200 has been debited in cash book but the double entry has not been completed.\par
}
\lqn{147}{In taking out a Trial Balance, a Bookkeeper finds that he is out Rs. 800 excess debit. Being desirous of closing his books, he places the difference to a newly opened suspense account. In the next period, he discovers the following discrepancies: Give the rectifying entries and prepare the suspense account.\\
(i) An item of sale for Rs. 5,900 was posted to the sale A/c as Rs. 9,500.\\
(ii) The total of the sale Returns Book has been added Rs. 100 short.\\
(iii) An amount of Rs. 3,700 received from a customer has been credited to his account as Rs. 7,300.\\
(iv) Rs. 1,50,000 paid for purchases of building have been charged to the ordinary purchase account.\\
(v) A sum of Rs. 9,500 written off from building A/c as depreciation has not been posted to depreciation A/c.\\
(vi) An amount received from a debtor of Rs. 9,700 has been debited to his account as Rs. 7,900.\par
}
\lqn{148}{A trade by preparing his Trial Balance found Rs. 780 excess debit and transferred it to the suspense account. Latter he found the following errors: (1) Pass the rectifying entries. (2) Prepare suspense account. (3) Ascertain correct profit after the rectification of errors-Before the rectification of errors the profit was Rs. 43870.\\
(a) Commission account balance (Dr.) Rs. 2,430 was not transferred to the Trial Balance.\\
(b) The total of the sales book Rs. 3,480 and not been posted to the sales A/c.\\
(c) Sale of old furniture for Rs. 2,400 was entered in the Sale Book.\\
(d) Purchases made from Raja Stores for Rs. 2,250 had been forgotten to enter in the books.\\
(e) Cash paid to Rehman Rs. 3,520 as posted to his account as Rs. 3,250.\\
(f) Goods sold to A for Rs. 1,250 posted to B’s Account.\par
}

\end{document}
